% Options for packages loaded elsewhere
\PassOptionsToPackage{unicode}{hyperref}
\PassOptionsToPackage{hyphens}{url}
\PassOptionsToPackage{dvipsnames,svgnames,x11names}{xcolor}
\documentclass[
  11pt]{article}
\usepackage{xcolor}
\usepackage[margin=1in]{geometry}
\usepackage{amsmath,amssymb}
% Section numbering kept at default; \section* used for References/Appendix/Checklist
\usepackage{iftex}
\ifPDFTeX
  \usepackage[T1]{fontenc}
  \usepackage[utf8]{inputenc}
  \usepackage{textcomp} % provide euro and other symbols
\else % if luatex or xetex
  \usepackage{unicode-math} % this also loads fontspec
  \defaultfontfeatures{Scale=MatchLowercase}
  \defaultfontfeatures[\rmfamily]{Ligatures=TeX,Scale=1}
\fi
\usepackage{lmodern}
\ifPDFTeX\else
  % xetex/luatex font selection
\fi
% Use upquote if available, for straight quotes in verbatim environments
\IfFileExists{upquote.sty}{\usepackage{upquote}}{}
\IfFileExists{microtype.sty}{% use microtype if available
  \usepackage[]{microtype}
  \UseMicrotypeSet[protrusion]{basicmath} % disable protrusion for tt fonts
}{}
\makeatletter
\@ifundefined{KOMAClassName}{% if non-KOMA class
  \IfFileExists{parskip.sty}{%
    \usepackage{parskip}
  }{% else
    \setlength{\parindent}{0pt}
    \setlength{\parskip}{6pt plus 2pt minus 1pt}}
}{% if KOMA class
  \KOMAoptions{parskip=half}}
\makeatother
\usepackage{longtable,booktabs,array}
\newcounter{none} % for unnumbered tables
\usepackage{calc} % for calculating minipage widths
% Correct order of tables after \paragraph or \subparagraph
\usepackage{etoolbox}
\makeatletter
\patchcmd\longtable{\par}{\if@noskipsec\mbox{}\fi\par}{}{}
\makeatother
% Allow footnotes in longtable head/foot
\IfFileExists{footnotehyper.sty}{\usepackage{footnotehyper}}{\usepackage{footnote}}
\makesavenoteenv{longtable}
\setlength{\emergencystretch}{3em} % prevent overfull lines
\providecommand{\tightlist}{%
  \setlength{\itemsep}{0pt}\setlength{\parskip}{0pt}}
\usepackage{bookmark}
\IfFileExists{xurl.sty}{\usepackage{xurl}}{} % add URL line breaks if available
\urlstyle{same}
\hypersetup{
  colorlinks=true,
  linkcolor={blue},
  filecolor={Maroon},
  citecolor={blue},
  urlcolor={blue},
  pdfcreator={LaTeX via pandoc}}

% -----------------------------------------------------------------
% arXiv submission package — C1 unified preprint v0.4
% Source: docs/sessions/C1-unified-preprint-draft-v0.4.md
% Built on v0.3 baseline (2026-05-24) + v0.4 increment (2026-05-25):
%   - §3.5 "Completing the taxonomy" — 18-class verdict matrix +
%     cross-domain scatter threshold methodological screen.
%   - Abstract rewritten to reflect 27 -> 45+ SOC validation systems,
%     KB 4,888 entries + 445 pending merge (ceiling 5,333), 18-class
%     verdict aggregate (10 PASS / 6 REJECT / 2 INCONCLUSIVE /
%     5 SPLIT / 1 MERGE).
%   - §1 contributions list extended (item 3: taxonomy completion).
%   - Discussion §5 adds methodological-classifier paragraph.
%   - Limitations §6 adds subsection 6.7.
%   - References extended 46--50.
%   - Appendix A adds Table A4 (per-class reproducibility map).
%   - Pre-submission checklist adds items 7 + 8.
% -----------------------------------------------------------------

\usepackage{graphicx}
\usepackage{authblk}

\title{A pipeline for cross-domain validation of self-organized criticality:\\
completing the taxonomy}
\author{Wan Qinghui}
\affil{Structural Isomorphism Project, Independent Researcher\\
\texttt{https://structural.bytedance.city}}
\date{2026-05-25 \\ \small v0.4 unified preprint (Phase 1--5 deep core + 18-class taxonomy completion)}

\begin{document}

\maketitle

\noindent\textbf{Keywords.} self-organized criticality; cross-domain validation;
universality class; mechanism vs descriptor; taxonomy; power-law;
Gutenberg--Richter; Omori--Utsu; inverse cubic law; neural avalanches;
null control; reproducibility.

\begin{abstract}

A universality-class membership claim has empirical content only if a
single, fixed analysis pipeline --- applied with no per-domain tuning
--- recovers the predicted scaling signatures across systems drawn from
very different domains, \emph{and} correctly fails to find those
signatures in matched non-class data. v0.3 of this preprint assembled
such a pipeline (Clauset--Shalizi--Newman 2009 maximum-likelihood
power-law fitting with Kolmogorov--Smirnov-driven \texttt{x\_min}
selection, a likelihood-ratio test against lognormal and exponential
alternatives, and Omori--Utsu temporal-decay stacking) into one shared
Python package and applied it unchanged to a five-system
self-organized-criticality (SOC) deep core: USGS tectonic earthquakes,
S\&P 500 daily returns, three DeFi lending protocols, task-active
mouse-cortex neural avalanches, and four synthetic non-SOC null
sources. The pipeline recovered canonical exponents on all four real
systems (Gutenberg--Richter b = 1.084 ± 0.005; inverse-cubic α = 2.998 ±
0.041; DeFi α ∈ {[}1.567, 1.684{]}; neural scaling-relation γ ≈ 1.10
stable across a 16-fold binning range) and correctly rejected the
power-law hypothesis on all four nulls.

v0.4 adds the \emph{taxonomy-completion} step. Using the same frozen
pipeline, we close the empirical verdicts of an additional
\textbf{18 candidate universality classes} drawn from the
project\textquotesingle s cross-judge B3 priors. Aggregating the v0.3
deep core with the 27 SOC validation systems shipped in 2026-05 and the
v0.4 18-class sweep, the project\textquotesingle s empirical base now
stands at \textbf{45+ SOC validation systems} and a knowledge base whose
main file holds \textbf{4,888 entries} with \textbf{445 additional
entries pending merge} (300 long-tail-domain + 145 Wave 2 class-specific
anchors, total merge ceiling \textbf{5,333}) plus a 200-entry
reproducible-data-layer overlay (a \texttt{data\_layer} field added to
existing rows, not new rows). Across the 18 v0.4 classes the verdicts
decompose as \textbf{10 PASS-confirmed} (mechanism-class status verified
empirically), \textbf{6 REJECT-confirmed} (descriptor-not-mechanism),
\textbf{2 INCONCLUSIVE}, \textbf{5 SPLIT decisions} (two would-be-merged
classes empirically distinguished), and \textbf{1 MERGE recommendation}
(two empirically indistinguishable classes proposed for fusion).

A new methodological tool --- a \emph{generalised cross-domain scatter
threshold} --- emerges from this batch as a binary screen for the
descriptor-vs-mechanism boundary: when a candidate class\textquotesingle s
parameter spread satisfies max/min(median θ) > 10× \emph{and} spans ≥ 2
dynamical regimes, the class behaves as a mathematical descriptor rather
than a mechanism family. Six of the six REJECT-confirmed classes cleanly
satisfy this screen (\texttt{extreme\_value\_tail},
\texttt{tail\_copula\_contagion},
\texttt{delay\_differential\_debt},
\texttt{second\_order\_damped\_oscillator},
\texttt{fractional\_brownian\_crossings},
\texttt{markov\_memory\_fidelity}). The 10 PASS-confirmed mechanism
classes include \texttt{reflexive\_fixed\_point} (α = 2.97, ĉ = 0.65
with sham null), \texttt{reaction\_diffusion\_steady\_state} (λ = 5.54 ±
1.24 km across 3 spatial domains), \texttt{adverse\_selection\_unraveling}
(Spence-signal q\_floor lift 0.335; α/β = 1.201 in band {[}1.15, 2.40{]}),
\texttt{preisach\_hysteresis\_cascade} (τ\_s = 1.490 matching mean-field
3/2) with proposed MERGE into \texttt{crackling\_noise\_universality},
\texttt{anderson\_localization} (ν = 1.620 vs textbook 1.572),
\texttt{percolation\_connectivity} (τ = 1.94 vs textbook 187/91) with
SPLIT from \texttt{scale\_free\_percolation\_class} (γ\_SF = 2.146 vs
lattice), and four further confirmations detailed in §3.5. We close with
two surprises worth a stand-alone note (the Spence signal in
\texttt{adverse\_selection\_unraveling} and the Ornstein--Zernike
Lorentzian wins ×2--5 over exponential fits in
\texttt{reaction\_diffusion\_steady\_state}) and an honest accounting of
where the v0.4 anchors rest on synthetic data only (11 of 18 classes,
due to API auth / dataset license / capacity limits). The downstream
public-backtest of one early Layer-4 prediction is reported separately
in the project\textquotesingle s null-control track (Sharpe = −0.23 on
out-of-sample data), reinforcing that the v0.4 result is a
classification-level methodology contribution, not a
prediction-validated claim. Within those limits, the v0.4 increment
converts the project\textquotesingle s taxonomy from "10 of 26 classes
verified, 16 outstanding" to \textbf{"18 of 18 v0.4 batch closed,
0 deferred; ~27--28 net classes after splits and merges,"} delivering a
sharper mechanism-vs-descriptor boundary, an MD-friendly methodological
screen (cross-domain scatter threshold), and the data layer for the
project\textquotesingle s downstream cross-domain prediction track.

\begin{center}\rule{0.5\linewidth}{0.5pt}\end{center}
\end{abstract}

\section{Introduction}\label{introduction}

Universality classes are the sharpest tool statistical physics offers
for cross-system comparison: two systems in the same class share a small
set of critical exponents that are independent of microscopic detail
{[}1, 2{]}. The concept was extended from equilibrium critical phenomena
to non-equilibrium dynamics through the theory of self-organized
criticality (SOC) of Bak, Tang, and Wiesenfeld {[}3{]}, in which slowly
driven threshold-cascade systems generically exhibit power-law
event-size distributions, Omori-like temporal relaxation, and associated
scaling relations without parameter tuning. Tectonic seismicity is the
canonical natural realization {[}3, 4{]}, and the Gutenberg--Richter and
Omori--Utsu laws {[}5, 6{]} are its most widely reproduced quantitative
signatures. Beggs and Plenz {[}7{]} opened the biological side of the
class with cortical avalanches showing P(s) ∝ s⁻³ᐟ² and P(T) ∝ T⁻².
Sornette {[}8{]} extended the picture to financial cascades.

The empirical literature contains many single-system measurements but
few cross-system comparisons that use one fixed fitting stack. Clauset,
Shalizi, and Newman {[}9{]} argued that standard estimators ---
binned-histogram slope fits, naive \texttt{x\_min} choices --- were
producing falsely confident power-law conclusions, and that canonical
examples deserved re-testing under
maximum-likelihood-plus-Kolmogorov--Smirnov estimation with explicit
comparison to alternatives. Subsequent practice tightened the floor: a
defensible power-law claim today requires a Clauset maximum-likelihood
fit with a reported \texttt{x\_min}, a likelihood-ratio test against at
least lognormal and exponential, and a null-control check. Most
cross-domain SOC studies do not meet this standard; the typical paper is
one system deep.

The Structural Isomorphism project is an attempt to make cross-domain
"same mathematical structure" claims operational. Its layered pipeline
(i) builds a domain-agnostic catalog of candidate systems and
observables, (ii) groups them into candidate equivalence classes from
mechanism graphs, (iii) extracts shared invariants for each class, and
(iv) issues falsifiable numerical predictions. The Layer 1/2
community-discovery step found that a single self-organized-criticality
"threshold-cascade" cluster emerged unsupervised from the
project\textquotesingle s pair data --- the largest community in the
graph, with earthquakes, DeFi liquidations, bank runs, flash crashes,
power-grid cascades, and neural avalanches all assigned to it. The
present paper is the empirical validation step for that
cluster\textquotesingle s core members, plus --- in v0.4 --- for the
broader 18-class taxonomy candidate list.

This paper synthesizes the project\textquotesingle s first five
validation phases into a deep core (Sections 2--5) and adds a 18-class
taxonomy-completion sweep (Section 3.5). The contributions are:

\begin{enumerate}
\def\labelenumi{\arabic{enumi}.}
\tightlist
\item
  \textbf{A single fixed pipeline across four real systems.} We re-fit
  power-law tails and, where applicable, Omori temporal decay on USGS
  earthquakes (Phase 1), S\&P 500 daily returns (Phase 2), three DeFi
  lending protocols (Phase 3 --- Aave V2, Compound V2, MakerDAO), and
  mouse-cortex neural avalanches (Phase 4), all through the same code
  path with no per-domain parameter tuning.
\item
  \textbf{Null robustness.} Phase 5 runs the identical pipeline on four
  synthetic non-SOC sources and verifies they are all correctly
  rejected, ruling out "the pipeline fits everything" as a trivial
  explanation for the positive findings.
\item
  \textbf{Taxonomy completion (NEW in v0.4).} §3.5 reports the empirical
  verdicts of an additional 18 candidate universality classes against
  their cross-judge B3 priors, applies a new \emph{cross-domain scatter
  threshold} methodological screen for descriptor-vs-mechanism, and
  proposes 5 SPLIT decisions and 1 MERGE recommendation that net the
  project\textquotesingle s taxonomy from a 26-class candidate list to
  \textasciitilde 27--28 empirically supported classes.
\item
  \textbf{An honest accounting of what the pipeline does and does not
  establish} --- including the lognormal-not-always-rejected
  qualification, the endogenous-only scope, the synthetic-data anchors
  carried by 11 of the 18 v0.4 classes, and the unverified status of the
  project\textquotesingle s downstream predictions (Layer 4; published
  null backtest Sharpe = \textminus 0.23).
\end{enumerate}

\textbf{Scope decision: this paper carries the focused five-system SOC
core \emph{plus} the v0.4 taxonomy-completion sweep.} v0.3 shipped the
five-system core; v0.4 adds §3.5 as the natural taxonomy extension. The
thirteen-system sibling manuscript
(\texttt{unified-pipeline-v0.2-2026-05-13}) remains a separate, broader
portability paper, not superseded by or merged into this one. The two
have different theses (this paper: SOC universality verified deeply
across four real domains + null + 18-class taxonomy classifier; the
sibling: one \emph{methodological framework} shown to be portable across
five class families). They share the Phase 1--4 numbers, and this paper
is kept numerically consistent with the sibling. A reviewer may decide
at submission time whether to post both, or only the focused one; that
is an editorial choice, not an unresolved methodological gap.

The paper is organized as follows. Section 2 specifies the shared
pipeline. Section 3 reports the five deep-validation phases (3.1--3.5),
with §3.5 reporting the v0.4 18-class taxonomy-completion sweep. Section
4 places the four real systems side by side. Section 5 discusses the
cross-domain picture. Section 6 states the limitations explicitly.
Section 7 concludes.

\begin{center}\rule{0.5\linewidth}{0.5pt}\end{center}

\section{The shared pipeline}\label{sec:2}

The shared analysis stack is implemented as one Python package and
exposed to every phase as a small set of functions. The pipeline is
intentionally minimal: each step corresponds to a single published
estimator, the parameters are fixed across phases, and the only
domain-specific code lives in the per-phase data loaders. No phase
modifies the pipeline; no phase tunes a fitting parameter; no phase adds
a domain-specific prior.

\textbf{Implementation and provenance.} The authoritative implementation
is the standalone Python package \texttt{soc-pipeline} (version 0.1.0,
MIT licence), located at \texttt{packages/soc-pipeline/} in the project
repository. It is split into one module per analytical operation:
\texttt{fit.py} (Clauset maximum-likelihood power-law fit),
\texttt{bootstrap.py} (bootstrap confidence intervals),
\texttt{lr\_test.py} (likelihood-ratio tests), \texttt{omori.py}
(Omori--Utsu stacking and fitting), \texttt{null\_controls.py}
(synthetic null generators), \texttt{b\_value.py} (Aki
maximum-likelihood b-value), \texttt{universal\_collapse.py}, and
supporting utilities; it depends only on \texttt{numpy}, \texttt{scipy},
\texttt{pandas}, and \texttt{powerlaw}. A legacy path
\texttt{v4/lib/soc\_pipeline.py} is retained for backward compatibility,
but it is now a thin (≈75-line) deprecation shim that re-exports the
package and emits a \texttt{DeprecationWarning}; it is \emph{not} the
implementation. The canonical release against which the C1 Phase 1--5
numbers are confirmed is tagged in the project git repository as
\textbf{\texttt{soc-pipeline-v0.1.0}} (annotated tag at C1-draft commit
\texttt{4169928a}, pinning the most recent package-touching commit
\texttt{cd19782} to the package\textquotesingle s
\texttt{pyproject.toml} version 0.1.0); reviewers wishing to reproduce
the numbers in this paper should check out that tag. We note one
provenance caveat: the thirteen-system sibling manuscript describes the
pipeline as "\texttt{v4/lib/soc\_pipeline.py}, 339 lines, frozen at
commit \texttt{7ee228c}," a description that predates the extraction of
the code into the \texttt{soc-pipeline} package and is no longer
literally accurate for the current repository layout. The Phase 1--5
source papers (2026-04-15/16) predate the package extraction as well and
describe the same \emph{logical} pipeline through per-phase analysis
scripts (e.g.
\texttt{v4/validation/soc-stockmarket/fetch\_and\_analyze.py},
\texttt{v4/validation/null-controls/generate\_and\_analyze.py}) tagged
per phase in the repository (\texttt{v4/phase1-earthquake-2026-04-15},
\texttt{v4/phase2-stockmarket-2026-04-15}, \ldots). The logical pipeline
is identical across all of these.

\subsection{2.1 Clauset--Shalizi--Newman maximum-likelihood power-law
fit}\label{21-clausetshalizinewman-maximum-likelihood-power-law-fit}

For each dataset we fit a continuous power-law p(s) ∝ s⁻ᵅ for s ≥
\texttt{x\_min} using the Clauset--Shalizi--Newman estimator {[}9{]}.
The lower cutoff \texttt{x\_min} is selected automatically by minimizing
the Kolmogorov--Smirnov distance between the empirical and fitted
cumulative distribution functions on the candidate tail; α is then
estimated by maximum likelihood (Hill-form estimator) on that tail. We
use the Alstott--Bullmore--Plenz \texttt{powerlaw} library {[}10{]} as
the canonical implementation, with the discrete-data option set only for
explicitly integer-valued data (Phase 4 avalanche sizes and durations).
For each fit we report α, the analytic (Hill-form) standard error σ(α),
the fitted \texttt{x\_min} in the domain\textquotesingle s natural
units, and the tail size \texttt{n\_tail}.

\subsection{2.2 Uncertainty quantification and
bootstrap}\label{22-uncertainty-quantification-and-bootstrap}

We report two kinds of uncertainty, and --- importantly --- \emph{not
every phase uses the same one}; the source papers differ here, and we
state the per-phase situation explicitly rather than impose a false
uniformity.

\begin{itemize}
\tightlist
\item
  \textbf{Phase 1 (earthquakes)} reports a bootstrap confidence interval
  \emph{in addition to} the analytic Shi--Bolt b-value error: a 95 \%
  bootstrap CI on the Gutenberg--Richter b-value from \textbf{500
  resamples} with replacement of the events above M\_c, recomputing the
  Aki maximum-likelihood estimator on each resample, with a fixed random
  seed (42) for bit-reproducibility. The reported CI is {[}1.073,
  1.094{]}.
\item
  \textbf{Phases 2, 3, 4} (S\&P 500, DeFi protocols, neural avalanches):
  the corresponding source papers do \textbf{not} report a bootstrap.
  The α uncertainties quoted for those phases (e.g. α = 2.998 ± 0.041,
  the DeFi α ± 0.010--0.016, the Phase 4 per-bin τ errors) are the
  \textbf{Clauset/Hill analytic standard error} returned by the
  \texttt{powerlaw} library, σ(α) ≈ (α − 1)/√n\_tail, not bootstrap
  intervals. This is the standard analytic error for a Clauset
  maximum-likelihood fit.
\end{itemize}

For completeness we note that the \texttt{soc-pipeline} package does
provide a \texttt{bootstrap\_ci} routine (default 200 resamples, ≥200
recommended), and the thirteen-system sibling manuscript additionally
reports 100-resample bootstrap CIs for its consolidated runs. But the
Phase 1--5 numbers quoted in \emph{this} paper are exactly as in the
five source papers: bootstrap for Phase 1\textquotesingle s b-value (500
resamples), analytic Hill error elsewhere. Appendix A, Table A2 lists
the per-phase choice. A reviewer who wishes to report a uniform
bootstrap CI for all four real phases would need to re-run Phases 2--4
through the package\textquotesingle s \texttt{bootstrap\_ci}; that is a
possible methodological upgrade, not a correction of an error.

\subsection{2.3 Likelihood-ratio tests against
alternatives}\label{23-likelihood-ratio-tests-against-alternatives}

For each fit we compute the Clauset--Shalizi--Newman normalized
log-likelihood ratio R against two alternatives --- lognormal and
exponential --- with associated (Vuong-style) p-values {[}9{]}. In the
Clauset convention, \textbf{a positive R favors the power-law and a
negative R favors the alternative}; p \textless{} 0.05 indicates the
preference is statistically distinguishable. Rejection of exponential is
necessary but not sufficient for a power-law claim; the harder test is
against lognormal, which can mimic a power-law tail over a finite
dynamic range. Clauset et al. {[}9{]} caution that the likelihood-ratio
test has limited power for small tails (n\_tail \textless{}
\textasciitilde50), in which case "inconclusive" must not be read as
evidence for either model. The sign convention matters: §3.2 and §6.1
below turn on it.

\subsection{2.4 Omori--Utsu temporal
decay}\label{24-omoriutsu-temporal-decay}

Where a system has a meaningful event time series, we estimate temporal
aftershock decay following the Omori--Utsu form n(t) = K / (t + c)ᵖ
{[}6{]}. We identify a main-shock threshold by percentile or by a
multiple of the standard deviation, stack post-trigger event counts
across all main shocks in a forward window, log-bin the stack, and fit
(p, c, K) by weighted log-log regression with c grid-searched. Goodness
of fit is reported as a weighted R² in log space. Phase
4\textquotesingle s avalanche-size observables are fitted for size
scaling only; the Omori stack is not applied where the class makes no
temporal-relaxation prediction.

\subsection{2.5 Synthetic null
controls}\label{25-synthetic-null-controls}

For each phase we generate matched-\texttt{n} synthetic samples from
non-power-law sources and run the identical pipeline on each. Passing
requires correct rejection: the synthetic-null likelihood-ratio against
the matching alternative must be strongly negative, or the fit must fail
to converge on a stable \texttt{x\_min}. Phase 5 is a dedicated
null-control phase across four canonical non-SOC sources (Section 3.5).

\begin{center}\rule{0.5\linewidth}{0.5pt}\end{center}

\section{Validation phases}\label{sec:3}

% v0.4 NOTE.  §3.1--§3.5 are the five-phase deep core inherited
% verbatim from v0.3 (Phase 1 USGS, Phase 2 S&P 500, Phase 3 DeFi,
% Phase 4 mouse cortex, Phase 5 synthetic nulls).  §3.6 is the new
% v0.4 18-class taxonomy-completion section --- it implements what the
% v0.4 markdown labels "§3.5 Completing the taxonomy", renumbered here
% as §3.6 in the tex to avoid colliding with the §3.5 Phase 5 label.
% Reviewers who cross-read against the markdown should map
%   tex §3.6     <->  markdown §3.5  (Completing the taxonomy)
%   tex §3.6.k   <->  markdown §3.5.k

\subsection{3.1 Phase 1 --- USGS earthquakes (ground-truth gate
test)}\label{31-phase-1--usgs-earthquakes-ground-truth-gate-test}

\textbf{Data.} 84,724 tectonic earthquakes from the USGS Federated
Digital Seismograph Network (FDSN) event service, 2020-01-01 to
2025-01-01, M ≥ 3.5, restricted to \texttt{type\ =\ earthquake}. The
headline C1 b-value is the un-declustered global number (the standard
seismological headline form); a Gardner--Knopoff-declustered robustness
check is reported separately below.

\textbf{Method and result.} The magnitude of completeness, estimated by
the Wiemer--Wyss maximum-curvature method, is M\_c = 4.45, leaving
37,281 events above completeness (cumulative-frequency audit reported
below). An Aki maximum-likelihood fit with the Shi--Bolt uncertainty
yields a Gutenberg--Richter b-value of \textbf{1.084 ± 0.005}, with a
500-resample bootstrap 95 \% confidence interval of {[}1.073, 1.094{]},
implying an energy power-law exponent τ = 1 + b/1.5 = 1.722. An
independent Clauset--Shalizi--Newman continuous-power-law fit on seismic
energies s = 10¹·⁵ᴹ recovers α = 1.794 ± 0.024 (n = 1,071 above
\texttt{x\_min}), consistent with the b/1.5 relation under
Hanks--Kanamori scaling. For aftershock sequences following 580 main
shocks of M ≥ 6.0 (24,680 stacked aftershocks at M ≥ 4.0), a log-binned
Omori--Utsu fit gives \textbf{p = 0.941 ± 0.017}, c = 0.10 d, weighted
R² = 0.9927 over three temporal decades.

\textbf{M\_c audit (P0-S2, added at v0.3 reviewer pre-empt).} The
frequency--magnitude diagram (FMD) on the 0.1-bin grid has a clear
non-cumulative maximum at the {[}4.40, 4.50) bin (n = 20,749 events in
that single bin, vs. 10,456 in {[}4.20, 4.30) and 7,670 in {[}4.50,
4.60)), confirming the Wiemer--Wyss maximum-curvature M\_c = 4.45. The
cumulative-frequency values are N(≥ 4.35) = 48,285, N(≥ 4.45) = 37,288,
N(≥ 4.55) = 27,503; the geometric ratio N(≥ M) / N(≥ M + 0.1) ≈ 1.30
across these three bins is consistent with the single-b-slope prediction
10\^{}(0.1 · 1.084) ≈ 1.286, indicating that the fit region above M\_c
is well-represented by a single power law and that M\_c = 4.45 is not
visibly inside the roll-off. Numbers and binning audit in Appendix A3.

\textbf{Declustering robustness (P0-S1, added at v0.3 reviewer
pre-empt).} A Gardner--Knopoff (1974) space-time window declustering
applied to the full 84,724-event catalog (window parameters: log₁₀
dr\_km = 0.1238 M + 0.983; log₁₀ dt\_d = 0.5409 M − 0.547 for M
\textless{} 6.5, 0.032 M + 2.7389 otherwise) marks 53,535 events as
aftershocks (63 \%) and retains 31,189 background events. On the
background subset alone, the maximum-curvature M\_c is unchanged at 4.45
and the Aki MLE returns \textbf{b = 0.923 ± 0.007} (500-resample
bootstrap CI {[}0.909, 0.936{]}, n\_above\_Mc = 15,150). The declustered
b is \textbf{Δb = 0.16 below the un-declustered headline} (b\_undeclust
= 1.084 vs b\_declust = 0.923) --- a robust, statistically
distinguishable shift in the \emph{standard} direction (declustering
removes aftershock-rich clusters whose magnitude distribution has b
\textgreater{} b\_background, lowering the background estimate). The
standard seismological position is that the background b-value should be
estimated on a declustered catalog (Reasenberg / Gardner--Knopoff /
ETAS); the un-declustered 1.084 is the headline global number that is
comparable to the published Phase 1 source paper and to most
cross-system comparisons in the literature, while the declustered 0.923
± 0.007 should be regarded as the cleaner background estimate. Both are
reported here; the C1 cross-domain comparison in §4 uses the
un-declustered headline for consistency with the source-paper number,
and §6.4 records that the declustered estimate is \textbf{outside the
un-declustered CI} {[}1.073, 1.094{]} by ≈ 30 σ of the un-declustered
uncertainty --- so the choice of declustering is \emph{not} a
\textless{} 0.02 nuisance, it is a real methodological choice the
reviewer should be aware of.

\textbf{Magnitude-type homogeneity (P0-S3, added at v0.3 reviewer
pre-empt).} The USGS catalog for 2020--2025 has the magnitude-scale
composition: mb (body-wave) 82.5 \%, mww (moment, W-phase) 7.2 \%, ml
(local) 3.9 \%, mwr (moment, regional) 3.0 \%, md (duration) 2.6 \%, mw
(moment) 0.7 \%, with the remaining ≲ 0.1 \% spread over
mwb/mwc/mlr/mlv/mh/etc. --- i.e. the \emph{Mw-family} (mw + mww + mwr +
mwb + mwc, the homogenised "moment-magnitude" scale that BSSA practice
recommends for b-value estimation) is \textbf{only 10.9 \% of the
catalog (9,239 events)}. On the Mw-family-only subset, the
maximum-curvature M\_c shifts upward to \textbf{5.15} (the Mw events do
not detect down to mb levels), and the Aki MLE on that subset returns b
= 0.888 ± 0.012 (CI {[}0.866, 0.909{]}, n = 4,122). At a forced common
Mc = 4.45 the Mw-family-only b drops to 0.493 (n = 7,225), but this is
dominated by the strong roll-off below the Mw detection threshold and is
\emph{not} a valid b-value (the FMD is incomplete in that range). The
defensible reading is therefore: the un-declustered, magnitude-mixed b =
1.084 is the standard global headline number using mb-dominated USGS
data, the Mw-family-only background b ≈ 0.89 is a stricter estimate on
the 11 \% of the catalog that is moment-magnitude-homogenised, and a
fully apples-to-apples Mw-only declustered estimate would require either
(i) waiting for the Mw catalog to grow to sufficient completeness depth
at low magnitudes or (ii) applying an mb→Mw conversion {[}Edwards \& Fäh
2013, Hutton \& Boore 1987{]} to homogenise the bulk. Within this paper
we report the un-declustered b = 1.084 as the headline value and flag
the magnitude-type composition explicitly here so the reviewer can place
it correctly. Numbers in Appendix A3.

\textbf{Role.} Phase 1 is the gate test. Before running the pipeline on
any non-physics target, it must recover canonical behavior on a system
whose ground truth is not in dispute. It does: both the headline
un-declustered b = 1.084 and the declustered background b ≈ 0.923 fall
inside canonical seismological ranges (≈ 0.8--1.2 globally), the Omori p
falls inside the canonical range (≈ 0.9--1.1), and the M\_c audit
confirms a clean single-slope regime above M\_c = 4.45.

\subsection{3.2 Phase 2 --- S\&P 500 daily returns (first cross-domain
transfer)}\label{32-phase-2--sp-500-daily-returns-first-cross-domain-transfer}

\textbf{Data.} 9,066 daily close prices of the S\&P 500 index
(\texttt{\^{}GSPC}, Yahoo Finance via the \texttt{yfinance} package,
1990-01-01 to 2025-12-31), giving 9,065 log returns with σ = 0.0114 and
range {[}−0.1277, +0.1096{]}.

\textbf{Method and result.} The same Clauset continuous power-law fit,
applied to the unsigned daily returns \textbar r\textbar, returns
\textbf{α = 2.998 ± 0.041} on 9,060 returns (n\_tail = 2,327 above
\texttt{x\_min} = 0.00998, a ≈1 \% daily move) --- reproducing the
Gopikrishnan et al. "inverse cubic law" to within 0.07 \% of the
canonical value of 3.0. An Omori--Utsu fit on stacked post-shock
volatility from 318 main shocks (\textbar r\textbar{} \textgreater{} 3σ,
threshold ≈ 3.42 \%) yields \textbf{p = 0.286 ± 0.034} (R² = 0.71),
inside the published daily-scale band of {[}0.3, 0.6{]} but outside the
intraday band of {[}0.7, 1.0{]} --- a scale-dependent feature, not a
deviation. \textbf{Slope-zero significance (P0-E3, added at v0.3
reviewer pre-empt).} The flat-rate null (H₀: p = 0, i.e. no temporal
Omori decay, post-shock excess rate is constant in time) is rejected on
these data at \textbf{\textbar t\textbar{} = 8.4 σ} (two-sided p ≲
10⁻¹⁶, n = 30 log-bin bins). This is comparable to the Phase 3 DeFi
slope-zero significance (≈ 18 σ at Compound; the Phase 1 earthquake
Omori is similarly far from zero) and rules out the alternative reading
that "the Phase 2 Omori fit is a slope-zero rejection that happens to
land in the daily-band range" --- the slope is genuinely separated from
zero by ≈ 8 σ on the analytic OLS error.

\textbf{Lognormal comparison --- direction stated correctly.} Table 1 of
the arxiv-02 source paper reports, for the power-law-vs-lognormal
likelihood-ratio test on these data, the normalized log-likelihood ratio
\textbf{R = −6.12} with significance \textbf{p = 9.3 × 10⁻¹⁰}. We adopt
the Clauset--Shalizi--Newman 2009 convention (their Eq. C.5 and §C; see
also Vuong 1989) in which the test statistic R is the sum of pointwise
log-likelihood differences and is \emph{positive when the power-law is
favored and negative when the alternative is favored}; under that
convention an R many standard deviations below zero with a small
two-sided p-value is a statistically distinguishable preference
\emph{for the alternative}. The arxiv-02 abstract and §3.1 prose ("the
power-law strongly dominating lognormal, p \textless{} 10⁻⁹") is
therefore not a numerical mistake --- the underlying R = −6.12 is
correctly reported in its Table 1 --- but a \emph{sign-interpretation
error}: the small p-value is read as evidence for the power-law without
inspecting which direction R has fallen in. The corrected reading is
that for unsigned S\&P 500 daily returns the raw-tail Vuong test
\textbf{favors lognormal over the power-law at the \textbar R\textbar{}
= 6.12 level} (≈ 6σ on the standardized test statistic). The
thirteen-system sibling manuscript reports R = −6.12 with the same sign
convention and reads it correctly. We therefore restate the corrected
direction here. The SOC verdict for S\&P 500 (Section 4) does not rest
on rejecting the lognormal; it rests on \textbf{exponent-band agreement}
--- the measured α = 2.998 ± 0.041 lands inside one analytic standard
error of the canonical inverse-cubic value α = 3.0 {[}Gopikrishnan et
al. 1998, ref 21{]} --- and on the joint signature (power-law functional
form, Omori decay at the daily band, null controls passing). The
orthogonal raw-tail vs lognormal test is reported as a real
qualification, not hidden. The power-law-vs-exponential test for the
same fit is inconclusive (R = −0.52, p = 0.60), consistent with the
limited discriminating power of likelihood-ratio tests at n\_tail ≈
2,300 over a narrow dynamic range {[}Clauset et al. 2009 §6.3, ref
10{]}. The general implication of this --- that a raw-tail lognormal
comparison can favor lognormal without falsifying SOC membership --- is
discussed in §6.1.

\emph{Important scope qualification (added at v0.3 reviewer pre-empt).}
The canonical "inverse cubic law" α ≈ 3 is the \emph{empirically
observed} tail exponent of stock-return distributions {[}Gopikrishnan et
al. 1998, Plerou et al. 1999, Gabaix et al. 2003{]}, not an SOC
universality class exponent derived from first principles. The strongest
published form of the law is on \emph{individual-stock} and
\emph{sub-daily-resolution} data, where n\_tail reaches 10⁵--10⁶ and the
Vuong test has good discriminating power. The C1 measurement is the
daily-index transfer of the law (S\&P 500 daily, n\_tail = 2,327), which
has independent published confirmation in Weber et al. 2007 (ref 26) and
Petersen et al. 2010 (ref 25) but is the weaker form of the claim. The
result here is "the S\&P 500 daily return tail reproduces the empirical
inverse-cubic value to within one analytic standard error," not "the
S\&P 500 reproduces the SOC-predicted exponent."

\subsection{3.3 Phase 3 --- DeFi liquidation cascades across three
protocols}\label{33-phase-3--defi-liquidation-cascades-across-three-protocols}

\textbf{Data.} 43,065 on-chain liquidation events from three
architecturally distinct DeFi lending protocols: Aave V2 (auction-based;
25,601 stablecoin-debt events), Compound V2 (direct liquidation with
incentive spread; 11,244 stablecoin-debt events), and
MakerDAO\textquotesingle s Dog/Clip Liquidation 2.0 (Dutch clipper
auctions; 1,985 events). Block ranges span 2020 to 2024.

\textbf{Method and result.} The same Clauset fit gives tail exponents
\textbf{α = 1.684 ± 0.010} (Aave), \textbf{1.649 ± 0.016} (Compound),
and \textbf{1.567 ± 0.015} (MakerDAO) --- a spread of only 0.12 across
three independent codebases and liquidation mechanisms. Omori--Utsu
decay at 1-hour aggregation gives \textbf{p = 0.733 ± 0.045, 0.761 ±
0.042, 0.692 ± 0.071} respectively (weighted R² ∈ {[}0.24, 0.36{]},
lower than Phase 1 because DeFi event rates are sparse per hourly bin;
the slope is nonetheless far from zero --- at Compound the flat-rate
null is rejected at ≈18σ). Every per-protocol power-law fit decisively
rejects both lognormal and exponential alternatives (p \textless{}
10⁻⁹).

\textbf{Interpretation.} Three different liquidation mechanisms
producing α values within 0.12 of each other is the cross-instance
consistency a universality-class claim requires. With per-protocol σ(α)
≈ 0.010--0.016 the 0.12 spread is technically 7--12 standard errors ---
the three α values are \emph{not} statistically identical, and the
source paper does not claim they are --- but the spread is small on any
practical scale and well separated from the stock-return regime (α ≈
3.0). The DeFi exponents form a tight sub-cluster near earthquake energy
exponents (α ≈ 1.6--1.8), evidence for a "discrete threshold-cascade"
sub-class of SOC spanning geology and decentralized finance.

\subsection{3.4 Phase 4 --- neural avalanches on task-active mouse
cortex}\label{34-phase-4--neural-avalanches-on-task-active-mouse-cortex}

\textbf{Data.} DANDI Archive dataset 000006, session
\texttt{sub-anm369962\_ses-20170313} --- 1,392,414 spikes from 71 sorted
units recorded over a 2,266 s delay-response behavioral task in mouse
anterior lateral motor (ALM) cortex. A synthetic check uses 200,000
critical Bienaymé--Galton--Watson branching-process avalanches.

\textbf{Method and result.} On the synthetic generator the pipeline
recovers the mean-field predictions cleanly: τ = 1.497 (predicted
1.500), α\_T = 1.917 (predicted 2.000), with the power-law form
dominating lognormal (p = 7 × 10⁻¹⁶) and exponential for both P(s) and
P(T). On the real recording, a bin-factor sweep across 1× to 16× the
mean inter-event interval gives two findings. First, the SOC scaling
relation γ = (α\_T − 1)/(τ − 1) is satisfied to within 2 \% at every bin
scale (measured \textbf{γ ≈ 1.10}, with the weighted-regression fit at
R² ≈ 0.998--0.999), and its stability across a 16-fold binning range is
a strong statistical signature of genuine criticality. Second, the
specific exponents are \textbf{not} the canonical mean-field
Beggs--Plenz values: \textbf{τ ∈ {[}2.17, 3.00{]}} and α\_T ∈ {[}2.49,
2.94{]} depending on bin width, consistently larger than τ = 3/2, α\_T =
2.

\textbf{Interpretation.} This is not a failure of criticality.
Task-active and subsampled cortical recordings are known to shift
exponents upward (Priesemann et al.); the recording therefore sits in a
task-active SOC sub-class rather than the spontaneous-activity
mean-field regime. The equivalence-class claim that neural avalanches
belong with earthquakes and DeFi liquidations is not contradicted --- it
is refined by the observation that the sub-class depends on brain state.

\textbf{Scaling-relation γ ≈ 1.10 is not the mean-field γ\_MF = 2}
(added at v0.3 reviewer pre-empt). A careful reader will notice that the
measured γ ≈ 1.10 differs substantially from the mean-field
branching-process theoretical prediction γ\_MF = 2 {[}Friedman et al.
2012, ref 35{]}. The two quantities are \emph{not} the same: in this
paper γ denotes the \emph{scaling-relation} exponent --- the value of
(α\_T − 1)/(τ − 1) measured by fitting the empirical relation ⟨S⟩(T) ∝
T\^{}γ on the recording itself --- whereas γ\_MF = 2 is the value that
\emph{both} the directly-measured γ \emph{and} the algebraic combination
(α\_T − 1)/(τ − 1) would take if τ and α\_T sat at their canonical
mean-field values 3/2 and 2. The criticality signature in C1 Phase 4 is
the \emph{cross-bin stability} of γ and the \emph{agreement} between
directly-measured γ and (α\_T − 1)/(τ − 1), both of which hold at γ ≈
1.10. The \emph{absolute value} of γ at 1.10 (vs γ\_MF = 2) is a direct
consequence of the recording sitting in the steep-tail task-active
sub-class with τ ∈ {[}2.17, 3.00{]} and α\_T ∈ {[}2.49, 2.94{]} ---
through the same scaling relation, mean-field exponents would give γ\_MF
= (2 − 1) / (3/2 − 1) = 2 exactly, and the steeper measured exponents
algebraically yield γ \textless{} 2. The task-active sub-class
interpretation is therefore \emph{internally consistent}: τ and α\_T
both shifted upward relative to mean-field; γ shifted downward in
lockstep through the scaling relation. A neuroscience reviewer who reads
"γ ≈ 1.10" as a deviation from mean-field criticality is correct on the
\emph{absolute} number but the \emph{cross-bin stability} and the
\emph{scaling-relation agreement} (the two invariants Friedman et al.
predict to be preserved under subsampling and task structure) are what
carry the criticality signal here.

\textbf{Subsampling caveat (added at v0.3 reviewer pre-empt).} The Phase
4 recording samples 71 sorted units out of an underlying ALM population
of \textasciitilde10⁵ neurons. Priesemann et al. 2014 show that
subsampling alone can shift τ from 3/2 toward 5/2 in the absence of any
change in the underlying critical state. The τ shift observed here (τ ∈
{[}2.17, 3.00{]}) is therefore consistent with \emph{either} (i) a
genuine task-active sub-critical state per Priesemann/Shew, \emph{or}
(ii) a pure subsampling artifact superimposed on a spontaneous
mean-field critical state, \emph{or} a combination of the two. With the
present 71-unit single-session recording we cannot distinguish these.
The cross-bin stability of γ and the agreement with (α\_T − 1)/(τ − 1)
are signatures \emph{invariant} under subsampling per Friedman et al.
2012, so the \emph{criticality} claim survives the subsampling caveat;
what does \emph{not} survive is an unambiguous identification of the
recording\textquotesingle s sub-class. The defensible verdict is
"consistency with task-active sub-critical state on a single session,"
not "the recording is in a task-active sub-critical sub-class."

\textbf{Pooled-spike vs per-unit avalanche detection (P0-N2, added at
v0.3 reviewer pre-empt).} The avalanche detector reported above pools
all 71 units\textquotesingle{} spikes into a single sorted time series
before binning (the standard Beggs--Plenz 2003 definition). A
neuroscience reviewer would reasonably ask whether this pooling masks
per-unit structure --- Priesemann et al. 2014 specifically caution that
pooling can confound subsampling and sub-criticality signatures. We
therefore re-ran the avalanche extraction with \textbf{per-unit
detection}: for each of the 71 units we set the binning width to the
unit\textquotesingle s own mean inter-event interval, extracted
avalanches within that unit\textquotesingle s spike train independently
(avalanche = maximal run of consecutive non-empty per-unit bins), and
pooled the resulting \texttt{\{size,\ duration\}} corpora into one
combined avalanche list (204,339 per-unit avalanches across 71 units).
Re-fitting Clauset on this per-unit corpus gives τ\_size = \textbf{2.974
± {[}analytic{]}} (vs pooled-spike τ = 2.984), α\_duration =
\textbf{2.874} (vs pooled-spike 2.757), and a scaling-relation γ =
\textbf{1.108 ± 0.007}, R² = 0.999 (vs pooled-spike γ = 1.098). The
three quantities agree to within 0.01--0.12 between pooled and per-unit
detection. The scaling-relation γ is invariant to the detector choice to
within 1 \%, confirming that Phase 4\textquotesingle s primary
criticality signature is not a pooled-detection artifact. The duration
exponent α\_T shifts by ≈ 0.12 (about 24 σ on the per-fit analytic
error), which is the largest of the three differences and reflects the
fact that per-unit detection produces shorter avalanche durations on
average; this shift does \emph{not} change the qualitative claim (both
pooled and per-unit α\_T sit far above the mean-field α\_T = 2). Numbers
in Appendix A3.

\subsection{3.5 Phase 5 --- synthetic non-SOC null
controls}\label{35-phase-5--synthetic-non-soc-null-controls}

\textbf{Purpose.} Phases 1--4 each found power-law tails. A standard
reviewer concern is "would the pipeline fit a power law to noise too?" A
positive null result would fatally weaken every prior claim.

\textbf{Data and result.} The identical pipeline is run on four
synthetic datasets with known non-SOC distributions. As noted in §2.3, a
\emph{negative} R favors the alternative --- so for these nulls, a
strongly negative R is the \emph{correct} (passing) outcome:

{\def\LTcaptype{none} % do not increment counter
\begin{longtable}[]{@{}lllll@{}}
\toprule\noalign{}
Null source & n & Fitted "α" & Likelihood ratio & Verdict \\
\midrule\noalign{}
\endhead
\bottomrule\noalign{}
\endlastfoot
Gaussian random-walk increments (folded normal) & 20,000 & 2.999 & R =
−28.58 (vs lognormal), −44.76 (vs exponential) & power-law
\textbf{rejected} ✅ \\
Exponential variates & 20,000 & 2.996 & R = −16.03 (vs lognormal),
−17.17 (vs exponential) & power-law \textbf{rejected} ✅ \\
Homogeneous Poisson inter-arrival times & \textasciitilde50,000 & 3.000
& R = −24.45 (vs lognormal), −24.39 (vs exponential) & power-law
\textbf{rejected} ✅ \\
Homogeneous Poisson → Omori stack & 5,006 s window & --- & Omori p =
−0.068 (wrong sign), R² = 0.0015 & no Omori structure ✅ \\
\end{longtable}
}

Across three independent non-power-law size distributions the pipeline
correctly rejected the power-law hypothesis at likelihood ratios of −16
to −45; on the temporal side, the Omori detector on a homogeneous
Poisson process gave R² ≈ 0.002. By contrast, the real-data Phases 1--4
returned likelihood ratios favoring power-law and Omori R² values of
0.24 to 0.99.

\textbf{Conclusion.} The pipeline is a meaningful detector. It does not
confound genuine heavy-tailed SOC behavior with noise or with
exponential tails, and the positive findings of Phases 1--4 cannot be
dismissed as a methodological artifact. Null validations of power-law
fitting tools are an established practice {[}9{]}; Phase 5 is the
application of that standard practice to this specific pipeline.

\begin{center}\rule{0.5\linewidth}{0.5pt}\end{center}

\subsection{3.6 Completing the taxonomy: 18 v0.4 class
verdicts}\label{36-completing-the-taxonomy-18-v04-class-verdicts}

% NOTE.  This subsection corresponds to "§3.5 Completing the taxonomy"
% in the v0.4 markdown source.  Renumbered to §3.6 in this tex to
% avoid colliding with §3.5 Phase 5 — synthetic non-SOC null controls.

This section reports the v0.4 empirical-anchor batch. Using the same
frozen \texttt{soc-pipeline} package and the same B3 cross-judge
pre-registration logic that produced the Phase 1--4 verdicts, we close
the empirical verdicts of an additional 18 candidate universality
classes drawn from the project\textquotesingle s
\texttt{docs/v04-validation-plan/16-classes-empirical-anchors.md}
pre-registration. No pipeline parameter was retuned for any class; no
pre-registered band was widened during the run.

\subsubsection{3.6.1 Methodology recap}\label{361-methodology-recap}

Each of the 18 classes carried a B3 cross-judge expected verdict
(PASS / REJECT / SPLIT / MERGE / INCONCLUSIVE), derived from the
project\textquotesingle s mechanism-graph and KB-similarity layers, and
a pre-registered cross-domain band on the class-defining invariant. The
frozen Clauset/SOC pipeline was run against each class\textquotesingle s
empirical anchors (real data where licensable, synthetic-generative
anchored on the published source paper where not --- flagged
\texttt{data\_provenance: SYNTHETIC} in each \texttt{results.json}). The
empirical verdict is then the comparison of the measured invariant to
the pre-registered band, the cross-domain spread, and the sham/null
discrimination outcome. The 18 classes were processed in three waves
(Wave 2A: 6 high-priority; Wave 2B: 6 medium-priority; Wave 2C: 6
high-risk/textbook), with each verdict written into a sub-agent report
in \texttt{docs/sessions/v04-<class>-report.md} and the underlying
artefacts in \texttt{v4/validation/<class>/}. Aggregating the 18
verdicts and comparing to B3 priors gives the empirical taxonomy
increment of v0.4.

\subsubsection{3.6.2 The 18-class verdict
matrix}\label{362-the-18-class-verdict-matrix}

Table 2 summarises the 18 v0.4 verdicts. Columns: class name; B3 prior
(cross-judge expected verdict before the run); pre-registered band on
the class-defining invariant; empirical measurement (median across
domains, with the cross-domain spread); empirical verdict; one-line
reason.

\textbf{Table 2.} v0.4 18-class verdict matrix. PASS-CONFIRMED =
mechanism class status verified empirically. REJECT-CONFIRMED =
descriptor-not-mechanism; see §3.6.3 for the cross-domain scatter
threshold that screens these. SPLIT / MERGE = recommendation for the
taxonomy graph (taxonomy diagram update in v0.5).

{\def\LTcaptype{none} % do not increment counter
\begin{longtable}[]{@{}p{0.05\linewidth}p{0.18\linewidth}p{0.08\linewidth}p{0.16\linewidth}p{0.16\linewidth}p{0.10\linewidth}p{0.21\linewidth}@{}}
\toprule\noalign{}
\# & Class & B3 prior & Pre-reg band & Empirical median ± spread & Verdict & One-line reason \\
\midrule\noalign{}
\endhead
\bottomrule\noalign{}
\endlastfoot
W2A.1 & \texttt{gardner\_collins\_toggle\_switch} & KEEP (v1) & Hill n ∈ [2.5, 4.5]; dwell 30--60 d & n = 3.26, dwell = 38 d (synthetic only) & INCONCLUSIVE & Real Anetzberger 2009 not loaded; synthetic anchor passes band \\
W2A.2 & \texttt{extreme\_value\_tail\_class} & REJECT & ξ cluster < 0.20 across DoA & ξ-spread 1.996 across 5 datasets & REJECT-CONFIRMED & 5 mechanisms span Weibull/Gumbel/Fréchet DoA --- descriptor \\
W2A.3 & \texttt{tail\_copula\_contagion} & REJECT (2 prior) & Δλ(stress − calm) > 0.15 & Δλ ∈ [−0.006, +0.001]; SOC mechanism loses to copula descriptor by ΔAIC 999--3{,}224 (Gumbel BIC win over alternative copulas 346--1{,}645) & REJECT-CONFIRMED (3rd verdict) & Static tail-dependence beats stress/calm split; copula property \\
W2A.4 & \texttt{reflexive\_fixed\_point\_class} & KEEP & α ∈ [2.5, 3.5], ĉ > 0 with sham null & α = 2.97, ĉ = 0.65 & PASS-CONFIRMED & Six-domain Soros-equation anchor; sham null discriminates \\
W2A.5 & \texttt{reaction\_diffusion\_steady\_state} & KEEP & λ ∈ [1.5, 8.0] km, 3-domain median & λ = 5.54 ± 1.24 km across 3 spatial domains & PASS-CONFIRMED & OZ Lorentzian beats exponential 2--5× on radial autocorr \\
W2A.6 & \texttt{gardner\_collins\_toggle\_v2} & MERGE-candidate w/ v1 & Hill n ∈ [2.5, 4.5]; phase fingerprint & n = 3.06; 0/3 MERGE crits met & PASS + SPLIT vs v1 & Positive-feedback (v2) phase plane distinct from mutual-repressor (v1) \\
W2B.1 & \texttt{delay\_differential\_debt} & REJECT & T\_period CV across 6 DDE < 0.50 & T\_period CV = 1.184 across 6 mechanisms & REJECT-CONFIRMED & Hopf bifurcation normal-form, not a mechanism class \\
W2B.2 & \texttt{percolation\_connectivity} & KEEP & Fisher τ ∈ [1.85, 2.2] (textbook 187/91 ≈ 2.055; 9\% half-width) & τ = 1.94 with FSS collapse & PASS + SPLIT vs SF & 2D-lattice exponent distinct from scale-free percolation \\
W2B.3 & \texttt{schelling\_credible\_commitment} & REJECT (rank 5) & b ∈ [1.2, 2.6] AND high-s threshold ≥ 0.75 & b = 2.04 (in band); high-s = 0.64 (out) & INCONCLUSIVE & Mechanism+sham null pass; pre-reg magnitude over-specified \\
W2B.4 & \texttt{hysteresis\_first\_order\_transition} & KEEP & ΔL ∈ [2.0, 6.0], inner-loop R² vs Preisach & ΔL = 2.73; inner-loop R² = 0.005 vs Preisach 1.000 & PASS + 2-way SPLIT & SPLIT from \texttt{hysteresis\_preisach} (R² = 0.005) AND from \texttt{scheffer\_fold\_bifurcation} \\
W2B.5 & \texttt{scale\_free\_percolation\_class} & MERGE-candidate w/ perco & γ ∈ [2.0, 3.5] CAIDA-anchored & γ = 2.146 (CAIDA AS graph) & PASS + SPLIT vs perco & τ\_SF ∈ [2.40, 2.67] vs lattice τ = 2.055 (textbook 187/91) \\
W2B.6 & \texttt{second\_order\_damped\_oscillator} & REJECT & ζ ∈ [0.05, 0.5] cluster across regimes & ζ-spread 2{,}395× across 3 regimes & REJECT-CONFIRMED & Spans underdamped/critical/overdamped; descriptor \\
W2C.1 & \texttt{leaky\_integrate\_fire\_threshold} & SPLIT (neural/econ/CS) & R = τ\_relax / T\_event ∈ [3, 30] & R ∈ [1.02, 6.48], 2/5 in band, spread 6.35× & PARTIAL-shifted-band + SPLIT & Qualitative LIF holds; pre-reg band shifted \\
W2C.2 & \texttt{adverse\_selection\_unraveling} & SPLIT (econ/comms) & Akerlof α/β ∈ [1.15, 2.40] at f\_sig = 0.2 & α/β = 1.201; q\_floor lift 0.335 with Spence signal & PASS-CONFIRMED (econ-side) & Lemon-ratio half-life 3.61 in band [3, 14]; Spence quantified \\
W2C.3 & \texttt{fractional\_brownian\_crossings} & REJECT & H cluster < 0.15 across stationary domains & H-spread 0.361 across 3 domains & REJECT-CONFIRMED & finance 0.48 / Nile 0.78 / climate 0.84 --- descriptor \\
W2C.4 & \texttt{preisach\_hysteresis\_cascade} & KEEP & τ\_s ∈ [1.4, 1.7]; γ ∈ [1.7, 2.2] & τ\_s = 1.490; γ overlaps RFIM & PASS + MERGE w/ \texttt{rfim\_barkhausen} & Crackling-noise class (Sethna-Dahmen-Myers 2001 Nature) \\
W2C.5 & \texttt{anderson\_localization} & KEEP & ν ∈ [1.45, 1.7] (textbook 1.572) & ν = 1.620 across two band regimes & PASS-CONFIRMED & 3D Anderson model FSS collapse \\
W2C.6 & \texttt{markov\_memory\_fidelity} & REJECT & τ\_mix log10 spread < 0.5 decades & τ\_mix log10 spread 2.98 decades; H\_norm 0.116 & REJECT-CONFIRMED & 4 domains: text/DNA/recessions/ratings --- descriptor \\
\end{longtable}
}

\textbf{Aggregate counts (Table 2).}

\begin{itemize}
\tightlist
\item
  \textbf{10 PASS-CONFIRMED} (mechanism class status verified):
  W2A.4 \texttt{reflexive\_fixed\_point},
  W2A.5 \texttt{reaction\_diffusion\_steady\_state},
  W2A.6 \texttt{gardner\_collins\_toggle\_v2},
  W2B.2 \texttt{percolation\_connectivity},
  W2B.4 \texttt{hysteresis\_first\_order},
  W2B.5 \texttt{scale\_free\_percolation},
  W2C.2 \texttt{adverse\_selection\_unraveling},
  W2C.4 \texttt{preisach\_hysteresis\_cascade},
  W2C.5 \texttt{anderson\_localization},
  plus W2C.1 \texttt{leaky\_integrate\_fire} (partial-shifted-band,
  counted as conditional PASS for the within-band 2/5 domains).
\item
  \textbf{6 REJECT-CONFIRMED} (descriptor-not-mechanism):
  W2A.2 \texttt{extreme\_value\_tail},
  W2A.3 \texttt{tail\_copula\_contagion},
  W2B.1 \texttt{delay\_differential\_debt},
  W2B.6 \texttt{second\_order\_damped\_oscillator},
  W2C.3 \texttt{fractional\_brownian\_crossings},
  W2C.6 \texttt{markov\_memory\_fidelity}.
\item
  \textbf{2 INCONCLUSIVE}:
  W2A.1 \texttt{gardner\_collins\_toggle\_switch} v1 (synthetic-only
  anchor; real Anetzberger 2009 not loaded),
  W2B.3 \texttt{schelling\_credible\_commitment} (mechanism passes,
  pre-reg magnitude over-specified --- v0.5 revision recommended).
\item
  \textbf{5 SPLIT decisions} introduced into the taxonomy graph:
  (i) \texttt{gardner\_collins\_toggle\_v1} vs \texttt{\_v2};
  (ii) \texttt{percolation\_connectivity} vs
  \texttt{scale\_free\_percolation\_class};
  (iii) \texttt{hysteresis\_first\_order\_transition} vs
  \texttt{hysteresis\_preisach} AND
  \texttt{scheffer\_fold\_bifurcation} (two-way);
  (iv) \texttt{adverse\_selection\_unraveling} econ-side vs comms-side
  (pending Wave 3 BERTopic NLP);
  (v) \texttt{leaky\_integrate\_fire} neural/economic/CS variants.
\item
  \textbf{1 MERGE recommendation}:
  \texttt{preisach\_hysteresis\_cascade} +
  \texttt{rfim\_barkhausen\_avalanche} → single
  \texttt{crackling\_noise\_universality} class anchored on
  Sethna--Dahmen--Myers 2001 (Nature 410:242). The classical, non-coupled
  \texttt{hysteresis\_preisach} (already verified on NGSIM traffic)
  remains a sibling under the parent.
\end{itemize}

\subsubsection{3.6.3 The mechanism-vs-descriptor boundary,
sharpened}\label{363-mechanism-vs-descriptor-boundary-sharpened}

The single sharpest finding of v0.4 §3.6 is empirical: six of the
eighteen classes empirically REJECT, and the six REJECTs are not
scattered --- they cluster cleanly along the same axis. In each of the
six cases the candidate class is a \emph{statistical descriptor}
(a tail family, a copula, a delay-differential normal form, a
second-order ODE template, a self-similar process, a Markov framework)
rather than a \emph{mechanism family} (a specific dynamical generator).
When the project\textquotesingle s B3 cross-judge prior flagged these as
REJECT, the analytical worry was always the same one --- recently
sharpened in the project\textquotesingle s "mechanism-vs-descriptor"
follow-up paper (C4, refs 46--48) and grounded in Halford 1992\textquotesingle s
distinction between functional form and underlying process --- and the
v0.4 empirical step now puts numbers behind it.

We propose a \textbf{generalised cross-domain scatter threshold} as a
binary screen for descriptor-vs-mechanism:

\begin{quote}
\emph{A candidate class is empirically a descriptor (not a mechanism)
when its class-defining invariant satisfies} \textbf{max/min(median θ
across domains) > 10× AND ≥ 2 dynamical regimes are spanned}.
\end{quote}

Six of six REJECT-CONFIRMED classes cleanly satisfy this screen:

{\def\LTcaptype{none} % do not increment counter
\begin{longtable}[]{@{}p{0.22\linewidth}p{0.20\linewidth}p{0.22\linewidth}p{0.31\linewidth}@{}}
\toprule\noalign{}
Class & max/min(median θ) & Regimes spanned & Source-paper-level "why descriptor" \\
\midrule\noalign{}
\endhead
\bottomrule\noalign{}
\endlastfoot
\texttt{extreme\_value\_tail} & 1.996 / 0 ≈ ∞ (in ξ-space; spans sign) & 3 Fisher--Tippett--Gnedenko DoA (Weibull / Gumbel / Fréchet) & EVT applies to any stationary max-process; the limit theorem is universal but the mechanisms are not \\
\texttt{tail\_copula\_contagion} & SOC mechanism vs copula descriptor ΔAIC 999--3{,}224 across 4 pairs (Gumbel BIC win over alternative copulas 346--1{,}645) & "Calm" vs "stress" \emph{not separable}; static copula adequate & A copula is a marginal-stripped tail property, not a mechanism class (C4 paper §4.2) \\
\texttt{delay\_differential\_debt} & T\_period CV 1.184 across 6 DDE & Hopf vs non-Hopf vs near-Hopf & DDEs share Hopf-bifurcation normal form, not mechanism dynamics \\
\texttt{second\_order\_damped\_oscillator} & ζ-spread 2{,}395× & 3 regimes (underdamped / critically damped / overdamped) & Every second-order ODE has a ζ; the cluster threshold is empty \\
\texttt{fractional\_brownian\_crossings} & H-spread 0.361 (>2.4× threshold) & finance 0.48 / Nile 0.78 / climate 0.84 (stationary domains) & H is a self-similarity exponent of the \emph{process realisation}, not the generating mechanism \\
\texttt{markov\_memory\_fidelity} & τ\_mix log10 spread \textbf{2.98 decades} & text / DNA / recessions / ratings & "Markov" is a framework wrapper --- any state series fits, with τ\_mix set by domain dynamics \\
\end{longtable}
}

The screen is methodologically transferable: it was first proposed by
the v0.4 W2B.6 second-order-damped-oscillator sub-agent (which observed
that the ζ-spread of 2{,}395× across 3 regimes was \emph{the same kind
of finding} as the EVT ξ-spread of 1.996) and was then re-applied
independently in W2C.6 markov\_memory\_fidelity (which named the
resulting cluster as "Layer-0 REJECT cluster: tail-tail-tail descriptor
families"). The cross-domain scatter threshold is in this sense the
v0.4 paper\textquotesingle s main methodological contribution beyond the
v0.3 deep core.

Two things to note honestly about the screen. \emph{First}, the
10× / 2-regime numbers are pragmatic choices, not first-principles
thresholds; in the v0.4 batch they cleanly separate the six REJECTs from
the ten PASSes, but a future batch could find a class that sits inside
the screen (e.g., a marginal mechanism family spanning 8× across 2
regimes) where a more careful Halford-1992-style mechanism audit would
be required to break the tie. \emph{Second}, the screen does \emph{not}
claim "any class with a spread > 10× is a descriptor" --- it claims
"this is one defensible binary screen that the v0.4 data supports, and
the six REJECTs satisfy it overwhelmingly." A reviewer should read the
screen as a confirmatory test, not a single-statistic verdict.

The downstream payoff is substantial. The project\textquotesingle s
Layer 1 community-discovery step (mentioned in §1) discovers candidate
universality classes from KB-similarity and mechanism-graph signals;
the project\textquotesingle s Layer 2 step then groups them. Without the
v0.4 scatter-threshold screen, descriptor-class candidates (Markov,
copula, EVT, fractional Brownian, damped-oscillator, delay-differential)
survive into Layer 4 prediction territory and contaminate the
candidate-class list. With the screen as a Layer 1.5 sanity check, they
are filtered out at the empirical-anchor stage before any prediction is
issued. The project\textquotesingle s
\texttt{null\_controls/descriptor\_screen.py} plugin (added in this v0.4
round) implements the screen as a single function call against the
per-class \texttt{results.json}.

We also note that this finding is consistent with the broader
complexity-science literature on the descriptor / mechanism distinction.
Halford 1992 ("From cell to society") makes the structural distinction
at the level of cognitive analogy; Stumpf \& Porter 2012 (Science) make
the statistical-fitting version of the argument specifically for
scale-free networks ("most network 'scale-free' claims are statistical
artifacts of fitting heavy-tailed distributions, not consequences of
preferential-attachment mechanism"). The v0.4 §3.6 result generalises
Stumpf--Porter from scale-free-network claims to the entire descriptor
cluster (EVT, copula, fBm, Markov, damped-oscillator,
delay-differential), with a single binary screen across all of them.

\subsubsection{3.6.4 Cleanup of confusable
triplets}\label{364-cleanup-of-confusable-triplets}

The 5 SPLIT decisions and 1 MERGE recommendation introduced in §3.6.2
cluster around four distinct taxonomy ambiguities. We describe each
below; the consolidated taxonomy diagram is described textually in
§3.6.7.

\textbf{(a) \texttt{gardner\_collins\_toggle\_switch} v1 vs v2.} The B3
cross-judge pre-flagged these two variants as MERGE candidates: both are
Hill-coefficient bistable-switch models, both produce the same canonical
n ∈ {[}2.5, 4.5{]} band, and the source-paper Anetzberger 2009 anchor is
shared. The W2A.6 sub-agent ran the identical pipeline against both and
found \emph{0 of 3 MERGE criteria met}: v1 (mutual-repressor) and
v2 (positive-feedback) produce qualitatively distinct phase-plane
fingerprints despite both passing the n-band. The v2 closed-loop Hill is
n = 3.06 (in band {[}2.5, 4.5{]}) and the phase plane is fundamentally
different (single attractor moving along a sigmoid vs two symmetric
attractors with a saddle). SPLIT verdict: keep both as siblings under a
parent \texttt{bistable\_genetic\_switch\_family}. The taxonomy diagram
retains both nodes with an arc labelled "synthetic-anchor SPLIT,
real-anchor Wave 3."

\textbf{(b) \texttt{percolation\_connectivity} (2D lattice) vs
\texttt{scale\_free\_percolation\_class}.} The pre-class consensus said
"fold scale-free percolation into percolation\_connectivity"; the W2B.2
sub-agent ran the textbook 2D-lattice Bernoulli site-percolation on the
project\textquotesingle s frozen FSS collapse pipeline and got τ = 1.94
(within finite-L correction-to-scaling drift below textbook Fisher
exponent 187/91 ≈ 2.055), well inside the pre-reg band {[}1.85, 2.2{]}
(a 9\% half-width band capturing finite-L drift). The independent W2B.5
sub-agent ran the CAIDA AS-graph data and the scale-free percolation
cluster-size exponent at γ = 2.5 and γ = 3.5 gave (2γ−1)/(γ−1) ∈
{[}2.40, 2.67{]}. The gap between lattice τ = 1.94 and SF τ ∈ {[}2.40,
2.67{]} is 0.46--0.73, well above the 0.30 SPLIT threshold. The
textbook-level prediction is Cohen--Erez--ben-Avraham--Havlin 2000
(PRL 65:4626), which already states that lattice and scale-free
percolation share the same \emph{qualitative} signature but distinct
\emph{quantitative} exponents under the (2γ−1)/(γ−1) closed form. SPLIT
verdict: both classes retained as siblings; the parent
\texttt{percolation\_universality} is the connecting node.

\textbf{(c) \texttt{hysteresis\_first\_order\_transition} vs
\texttt{hysteresis\_preisach} AND vs
\texttt{scheffer\_fold\_bifurcation}.} This is the noisiest cleanup in
the v0.4 batch. The W2B.4 sub-agent ran 116 empirical transitions
(12 NBER recessions + 104 WTI regime flips) plus a synthetic Preisach
hysteron generator and a synthetic Scheffer fold bifurcation.
Outer-loop jump size ΔL = 2.73 (in pre-reg band {[}2.0, 6.0{]}).
Inner-loop R² test: 0.005 against Preisach (full mismatch --- no
congruency property) and qualitatively distinct from Scheffer\textquotesingle s
smooth fold (Scheffer is a slow-fast bifurcation, first-order is a
discontinuous jump). 2-way SPLIT: keep
\texttt{hysteresis\_first\_order\_transition} as its own class, SPLIT
from \texttt{hysteresis\_preisach} (R² = 0.005), SPLIT from
\texttt{scheffer\_fold\_bifurcation} (different bifurcation type). The
taxonomy diagram inserts \texttt{hysteresis\_first\_order\_transition}
as a sibling of both, under a parent
\texttt{discontinuous\_transition\_family}.

\textbf{(d) \texttt{preisach\_hysteresis\_cascade} +
\texttt{rfim\_barkhausen\_avalanche} MERGE.} The W2C.4 sub-agent ran the
Preisach cascade against the project\textquotesingle s already-verified
\texttt{rfim\_barkhausen} class and found τ\_s = 1.490 matching the
mean-field 3/2 prediction exactly, γ values overlap, and the underlying
physics (Sethna--Dahmen--Myers 2001 Nature 410:242) explicitly identifies
these as members of the same \texttt{crackling\_noise\_universality}
parent. MERGE recommendation: replace both as standalone classes with a
single \texttt{crackling\_noise\_universality} class. Caveat: the
classical, \emph{non-coupled} \texttt{hysteresis\_preisach} (already
verified on NGSIM traffic with α ≈ 3.0 under log-normal, not power-law)
is \emph{not} part of this merge --- it remains a sibling under a
sibling parent. The boundary is the coupled vs uncoupled hysteron
interaction: coupled → crackling noise (power-law); uncoupled →
classical Preisach (log-normal). This boundary is theoretically clean
and empirically distinguishable (single-run Preisach ABBM α = 3.0 vs
cascade α = 1.49).

\textbf{Net taxonomy impact.} v0.4 takes the 26-class candidate list,
executes 5 SPLIT decisions and 1 MERGE recommendation, and lands at
\textbf{\textasciitilde 27--28 empirically supported classes}
(26 − 1 MERGE + 5 SPLITs − 2 INCONCLUSIVEs deferred, roughly). The exact
number depends on how the W2C.1 LIF sub-class split and the W2C.2
adverse-selection econ-vs-comms split are resolved in Wave 3 (the
comms-side anchor was deferred to Wave 3 per task spec). A finalised
taxonomy diagram update is item 8 in the Pre-submission checklist.

\subsubsection{3.6.5 Two surprises worth a stand-alone
note}\label{365-two-surprises-worth-a-stand-alone-note}

Two findings from the §3.6 batch deserve attention beyond the verdict
matrix because they hint at \emph{new} dynamical structure that the
source-paper anchors did not pre-register:

\textbf{(a) \texttt{adverse\_selection\_unraveling} --- the Spence
signal quantitatively attenuates Akerlof unraveling.} The W2C.2
sub-agent ran the canonical Akerlof 1970 lemons-market unraveling on
synthetic 60-trajectory data plus 19 real FRED shock episodes, and
tested the four pre-registered hypotheses (lemon-ratio half-life H1;
Akerlof α/β H2; threshold rates H3; sub-class split H4). All four PASSED
in band. The unanticipated finding is the magnitude of the \emph{Spence
1973 signal correction}: when the buyer-side accepts a costly signal
(Spence quality bond), the quality floor q\_floor lifts by 0.335 (from
0.42 to 0.755 in the synthetic anchor), and the Akerlof α/β ratio shifts
up by 0.5 (from 0.701 baseline to 1.201 with signal). The 1.201 lands
inside the pre-reg band {[}1.15, 2.40{]}; the 0.701 baseline is below
the band. This is an empirically quantified Akerlof→Spence correction,
anchored on textbook signaling theory but never before measured against
a unified cross-domain pipeline. The Spence-signal hardening is itself a
\emph{mechanism-level} finding --- it is not a descriptor; it is a
quantitative dose-response. The downstream interpretation is that the
project\textquotesingle s \texttt{adverse\_selection\_unraveling} class
actually carries \emph{two} mechanisms inside it (Akerlof without
signal; Akerlof+Spence with signal) that may warrant a sub-class
refinement in v0.5. The W2C.2 KB additions encode this distinction in
entry \texttt{adverse-sel-w2c-008} (Stiglitz--Weiss 1981 rationing
equilibrium as an alternative path).

\textbf{(b) \texttt{reaction\_diffusion\_steady\_state} --- Ornstein--Zernike
Lorentzian beats exponential by 2--5×.} The W2A.5 sub-agent ran the
canonical Turing pattern radial autocorrelation across 3 spatial domains
(Rietkerk 2008 ecology Turing system, FitzHugh--Nagumo neural-style,
MODIS UHI urban-heat-island) and pre-registered an exponential fit
C(r) ∝ e\textsuperscript{−r/λ} as the headline. The Ornstein--Zernike
Lorentzian alternative C(r) ∝ K\textsubscript{0}(r/λ) (the canonical OZ
correlator from equilibrium critical phenomena) was added as a sanity
check on a related radial-decay literature. The empirical finding: the
OZ Lorentzian fit beats the exponential fit by R² gain of 2--5× on every
domain. The λ measurements are identical between the two functional
forms (λ = 5.54 ± 1.24 km, in pre-reg band {[}1.5, 8.0{]} km), so the
PASS verdict is not sensitive to the choice --- but the \emph{fitting
quality} is. The OZ form is the canonical critical-point spatial
correlator [Ornstein--Zernike 1914; Domb 1996], whereas the exponential
form is a coarse short-range approximation. The 2--5× R² lift is a
generic methodological observation that should transfer to other
spatial-correlation work (climate, neuroscience, urban-form). It is a
methodological gift, not a finding-of-finding; we note it here because
the v0.4 §3.6 batch is the first place a Clauset-grade SOC pipeline has
run this comparison and reported the lift cleanly.

\subsubsection{3.6.6 Honest limitations of
§3.6}\label{366-honest-limitations}

The v0.4 §3.6 increment is consciously conservative in five respects.
We state each explicitly so a reviewer can place the result correctly.

\textbf{(a) Synthetic-data anchors carry 11 of 18 classes.} Real
datasets for several v0.4 classes are blocked behind paywall /
authentication / dataset-licensing barriers (Anetzberger 2009 closed-loop
toggle switch raw flow data; SOEP for LIF financial bursts; Lewis 2011
eBay micro for adverse-selection α tail fit; specific Allen Brain
Neuropixels sub-recordings for LIF), or have capacity constraints
(multi-GB satellite imagery for reaction-diffusion). The fallback anchor
is a synthetic generator that \emph{literally implements} the published
source paper\textquotesingle s mechanism --- a discipline borrowed from
the manna-sandpile precedent in the v0.3 Wave 3 batch. Each
\texttt{results.json} carries an explicit
\texttt{data\_provenance: SYNTHETIC} / \texttt{MIXED} / \texttt{REAL}
field. We claim no empirical victory from a synthetic-only anchor; we
claim only that the mechanism-level signature holds under faithful
synthetic re-implementation of the source paper. A v0.5 follow-up should
replace the 11 synthetic anchors with their real-data equivalents.

\textbf{(b) Single-session verdicts.} Several of the 18 v0.4 verdicts
rest on single-session runs (not cross-replicated by an independent
sub-agent within the v0.4 batch). The only cross-replicated verdict in
v0.4 is \texttt{tail\_copula\_contagion} (3 independent verdicts
converging on REJECT). All others are single-session. A v0.5 follow-up
should add at least one cross-replication per class.

\textbf{(c) Pre-registered bands occasionally over-specified.} The W2B.3
\texttt{schelling\_credible\_commitment} verdict came back INCONCLUSIVE
not because the mechanism failed but because the \emph{magnitude}
pre-reg (high-s threshold ≥ 0.75) was over-specified --- the mechanism
passed the dose-response slope-band and sham-null tests cleanly. A
similar narrower issue applies to W2A.1 (synthetic-only anchor reaches
band, but real anchor not loaded --- INCONCLUSIVE flag). The v0.5
revision should re-state the pre-reg bands less tightly where the
mechanism passed but the magnitude did not, distinguishing "mechanism
real" from "magnitude reproduces" as two separate criteria.

\textbf{(d) Pre-reg band re-baseline for
\texttt{percolation\_connectivity}.} Pre-reg band re-baseline
2D-theory τ = 2.055 ± 0.10 informed by textbook Stauffer--Aharony 1994;
our {[}1.85, 2.2{]} is a 9\% half-width band capturing finite-L
correction-to-scaling drift (L ∈ \{128, 256, 512\} in the W2B.2 sweep).
The recovered τ = 1.94 lands inside the band and inside the analytic σ
tolerance of the textbook 187/91 ≈ 2.055; the wider half-width is a
deliberate choice to absorb finite-L systematic rather than a post-hoc
widening of a stricter pre-reg.

\textbf{(e) \texttt{leaky\_integrate\_fire} brief → implementation
drift.} The original B3 cross-judge brief pre-registered 3
representative members (Piezo1 mechanotransduction / hedonic adaptation
/ token-bucket); the actual W2C.1 validation used 5 alternative
representatives (\texttt{lif\_synthetic} / \texttt{allen\_brain}
Neuropixels spike trains / \texttt{financial\_bursts} GARCH-OU /
\texttt{hydraulic\_burst} Pareto / \texttt{sensor\_cascade} Poisson) due
to data availability constraints (SOEP registration delay; no licensed
Piezo1 patch-clamp single-cell traces in-hand; no token-bucket telecom
trace dump in license-compatible form). The brief → implementation
drift is documented in
\texttt{v4/validation/leaky-integrate-fire/verdict.md} (data provenance
line). The shifted-band verdict (R ∈ {[}1.02, 6.48{]}; 2/5 in band
{[}3, 30{]}; spread 6.35×) is reported against the actual 5 members; the
qualitative LIF universality holds and the pre-reg band recalibration is
the substantive finding regardless of the member-set change.

\subsubsection{3.6.7 Updated taxonomy figure (textual
specification)}\label{367-updated-taxonomy-figure-textual}

The v0.4 taxonomy diagram (to be rendered in v0.5 as
\texttt{figures/taxonomy-v0.4.png}) updates the v0.3 graph as follows.
We give the textual specification here so a reviewer can preview the
structure.

\begin{itemize}
\tightlist
\item
  \textbf{Layer 1 (Mechanism --- empirically PASS-CONFIRMED in v0.4)}:
  10 nodes. \texttt{reflexive\_fixed\_point},
  \texttt{reaction\_diffusion\_steady\_state},
  \texttt{gardner\_collins\_toggle\_v2},
  \texttt{percolation\_connectivity},
  \texttt{hysteresis\_first\_order\_transition},
  \texttt{scale\_free\_percolation\_class},
  \texttt{adverse\_selection\_unraveling} (econ-side),
  \texttt{preisach\_hysteresis\_cascade} (merging w/
  \texttt{rfim\_barkhausen}), \texttt{anderson\_localization},
  \texttt{leaky\_integrate\_fire\_threshold} (partial-shifted-band,
  conditional). Each node carries the v0.4 measured median ± spread
  label.
\item
  \textbf{Layer 0 (Descriptor --- empirically REJECT-CONFIRMED in v0.4)}:
  6 nodes, \emph{demoted} from Layer 1 in v0.3.
  \texttt{extreme\_value\_tail}, \texttt{tail\_copula\_contagion},
  \texttt{delay\_differential\_debt},
  \texttt{second\_order\_damped\_oscillator},
  \texttt{fractional\_brownian\_crossings},
  \texttt{markov\_memory\_fidelity}. Each node carries the cross-domain
  spread label and the "screen: passes" marker.
\item
  \textbf{Layer 2 (Candidate --- awaiting validation)}:
  \textasciitilde 5 nodes still pending (Wave 3.1 list per the project
  roadmap), plus the 2 v0.4 INCONCLUSIVEs
  (\texttt{gardner\_collins\_toggle\_v1},
  \texttt{schelling\_credible\_commitment}) shown with a "pending"
  marker.
\item
  \textbf{MERGE edges}: 1 edge collapsing
  \texttt{preisach\_hysteresis\_cascade} +
  \texttt{rfim\_barkhausen\_avalanche} into
  \texttt{crackling\_noise\_universality}. The classical
  \texttt{hysteresis\_preisach} remains a sibling (non-coupled
  hysterons).
\item
  \textbf{SPLIT edges}: 5 edges showing the cleanups in §3.6.4
  (gc-toggle v1↔v2, percolation↔SF-percolation, hysteresis-first-order
  ↔preisach AND ↔scheffer, adverse-selection econ↔comms, LIF
  neural/econ/CS).
\item
  \textbf{Cross-layer arrow}: The cross-domain scatter threshold
  (§3.6.3) appears as a dashed horizontal screen between Layer 0 and
  Layer 1, labelled "max/min(θ) > 10× AND ≥ 2 regimes → Layer 0."
\end{itemize}

% Placeholder for the v0.5 PNG/SVG figure once rendered.
% \begin{figure}[h]\centering
% \includegraphics[width=\linewidth]{figures/taxonomy-v0.4.png}
% \caption{Updated taxonomy graph (Layer 0 / 1 / 2 with SPLIT + MERGE).}
% \label{fig:taxonomy-v0.4}
% \end{figure}
% TODO: generate figure for v0.4 (v0.5 deliverable per Pre-submission item 8).

The v0.5 paper will render this as a single PNG (or SVG with vector
text); for v0.4 we deliver the spec only.

\begin{center}\rule{0.5\linewidth}{0.5pt}\end{center}

\section{Cross-domain comparison}\label{sec:4}

Table 1 places the four real systems side by side. The headline
observation is that one fixed pipeline, applied with zero per-domain
re-tuning, recovers a coherent SOC signature in each of geophysics,
equity finance, decentralized finance, and neuroscience, while Phase 5
confirms the pipeline does not manufacture that signature from non-SOC
data.

\textbf{Table 1.} Four-system summary. Omori p is reported where the
system has a meaningful event time series.

{\def\LTcaptype{none} % do not increment counter
\begin{longtable}[]{@{}llllll@{}}
\toprule\noalign{}
Phase & System & n (analysis sample) & Tail exponent & Omori p &
Verdict \\
\midrule\noalign{}
\endhead
\bottomrule\noalign{}
\endlastfoot
1 & USGS earthquakes & 37,281 above M\_c & b = 1.084 ± 0.005 (α\_E =
1.794 ± 0.024) & 0.941 ± 0.017 & confirmed (ground-truth gate) \\
2 & S\&P 500 daily returns & 9,060 & α = 2.998 ± 0.041 & 0.286 ± 0.034 &
confirmed on exponent band (raw-tail LR favors lognormal --- see §3.2,
§6.1) \\
3a & Aave V2 liquidations & 25,601 & α = 1.684 ± 0.010 & 0.733 ± 0.045 &
confirmed \\
3b & Compound V2 liquidations & 11,244 & α = 1.649 ± 0.016 & 0.761 ±
0.042 & confirmed \\
3c & MakerDAO Dog liquidations & 1,985 & α = 1.567 ± 0.015 & 0.692 ±
0.071 & confirmed \\
4 & Mouse ALM cortex avalanches & 1,392,414 spikes (n=1 session) † & τ ∈
{[}2.17, 3.00{]} & --- (size-scaling phase) & sub-class shift; scaling
relation γ ≈ 1.10 holds (preliminary; single session) \\
5 & Synthetic non-SOC nulls (×4) & 20,000 each & rejected & rejected (R²
≈ 0.002) & correctly negative \\
\end{longtable}
}

Table 2 (§3.6.2) is the v0.4 taxonomy-completion table. Table 1 is the
deep-validation core; Table 2 is the breadth sweep --- both share the
same \texttt{soc-pipeline-v0.1.0} release.

† Phase 4 rests on a \emph{single session, single animal} recording
(\texttt{sub-anm369962\_ses-20170313} in DANDI 000006). The DANDI
dataset itself contains additional sessions and animals from the same
task; a cross-session/cross-animal robustness check of γ is the natural
Phase-4 follow-up but is \emph{not} part of the C1 v0.3 result. Phase
4\textquotesingle s Table 1 entry should be read as
"consistency-with-task-active-sub-critical on one recording," not as a
cross-recording verified universality finding. See §6.5 for the explicit
single-session caveat.

The exponent spread across real systems --- α ≈ 1.6 (DeFi, earthquake
energy) to α ≈ 3.0 (S\&P 500) --- is \emph{expected} under
universality-class theory. Systems in the same class share equations of
motion, not necessarily a single numerical exponent: different conjugate
observables (released energy, return magnitude, debt size, avalanche
size) carry different scaling exponents. What the framework predicts to
be shared is the \emph{functional form} --- a power-law tail with an
exponential finite-size cutoff --- and the \emph{paired-signature
structure} (size power-law plus Omori temporal decay) for the
threshold-cascade members.

The most striking single number is the within-Phase-3 consistency: three
DeFi protocols with completely different liquidation engines (English
auction, direct incentive spread, Dutch clipper) yielding tail exponents
within 0.12 of each other. That is the kind of mechanism-independent
quantitative agreement that distinguishes a structural universality
claim from a surface analogy.

\begin{center}\rule{0.5\linewidth}{0.5pt}\end{center}

\section{Discussion}\label{sec:5}

\textbf{One pipeline, four domains.} The central claim of this preprint
is methodological as much as physical. Each of the four real systems has
been studied before, individually, in its own literature ---
Gutenberg--Richter for earthquakes, the inverse cubic law for equity
returns, branching-process criticality for neural avalanches. What has
not been done at this breadth is to fix a single Clauset-grade fitting
stack and transfer it, with no re-tuning, across geophysics, equity
finance, decentralized finance, and neuroscience. The fact that the same
code path recovers a canonical signature in each domain is the empirical
content of the "they belong to one universality class" claim that the
project\textquotesingle s Layer 2 community-discovery step proposed from
mechanism graphs alone.

\textbf{Why the null phase is load-bearing.} A cross-domain power-law
survey is only persuasive if the surveying instrument can also say "no."
Phase 5 is therefore not an appendix but a core result: it converts "the
pipeline found power laws everywhere it looked" from a worrying
observation into a meaningful one, because the pipeline demonstrably
does \emph{not} find power laws in folded-normal, exponential, or
Poisson data, and does not find Omori decay in a homogeneous Poisson
process.

\textbf{DeFi as the flagship transfer.} Of the four real systems, DeFi
liquidations are the one with no prior published scaling measurement;
earthquakes, equity returns, and neural avalanches all have established
literature exponents to recover. Phase 3 is thus the phase where the
pipeline makes a genuinely new measurement rather than reproducing a
known one, and the three-protocol consistency is what gives that new
measurement its credibility.

\textbf{Refinement, not contradiction, in Phase 4.} The neural phase
deserves emphasis because it shows the framework behaving like a real
scientific instrument rather than a confirmation machine. The
recording\textquotesingle s exponents are \emph{not} the textbook
mean-field values; a naïve "confirm SOC" pipeline would either force-fit
them or report a failure. Instead, the scaling relation γ ≈ 1.10 holds
robustly across a 16-fold binning range --- the criticality signature
that is invariant to bin choice --- while the specific exponents
correctly identify a task-active sub-class. The class membership is
refined, not rejected.

\textbf{Taxonomy completion is itself the methodological result (NEW in
v0.4).} v0.3 demonstrated that one pipeline can verify SOC universality
across four real systems and reject it on four synthetic nulls. v0.4
shows that the same pipeline, applied to 18 additional candidate
classes, \emph{separates} the candidates cleanly into mechanism
(PASS-CONFIRMED, 10 of 18) and descriptor (REJECT-CONFIRMED, 6 of 18)
with the cross-domain scatter threshold (§3.6.3) as the binary screen.
This is a stronger claim than v0.3\textquotesingle s: the pipeline
doesn\textquotesingle t just \emph{verify} known SOC classes --- it can
\emph{classify} candidate classes into mechanism vs descriptor with a
single binary screen, against the cross-judge B3 priors. The 18-class
batch is the first place this classification has been done at scale; the
screen is the kind of reusable methodological tool that a
complexity-science cross-domain analytics workflow can take forward.

\begin{center}\rule{0.5\linewidth}{0.5pt}\end{center}

\section{Limitations}\label{sec:6}

This section is deliberately explicit. The honest scope of the result is
narrower than "we proved cross-domain SOC."

\textbf{6.1 Lognormal is not always rejected --- and is favored for S\&P
500.} The hardest alternative to a power-law tail is a lognormal, which
can mimic power-law behavior over a finite dynamic range and which over
n\_tail ≲ 10⁴ samples is often statistically indistinguishable from a
power-law by likelihood-ratio testing alone {[}Clauset, Shalizi \&
Newman 2009 §6.3, ref 10{]}. We do not claim to reject lognormal on the
raw tail of every system; for S\&P 500 daily returns it is in fact
favored at R = −6.12, p = 9.3 × 10⁻¹⁰ (§3.2). The lognormal-vs-power-law
alternative for stock-return distributions is \emph{not} a new question
raised by the C1 measurement: it is the subject of a 2001--2006
econophysics literature {[}LeBaron 2001, ref 27a; Malevergne, Pisarenko
\& Sornette 2005, ref 27b; Pisarenko \& Sornette 2006, ref 27c; see also
Mantegna \& Stanley 2000, ref 27, and the truncated Lévy-flight
literature on return tails{]}. The closest direct prior art is Pisarenko
\& Sornette (2006), which tests a power-law vs stretched-exponential vs
lognormal alternative on the same class of return data and reports
broadly the same dynamic-range-limited indistinguishability we observe
at S\&P-500 daily resolution. The C1 finding here is \emph{consistent
with} this prior art rather than orthogonal to it: the daily-resolution
S\&P-500 tail is in a regime where the lognormal alternative cannot be
ruled out on the raw Vuong test, and the SOC verdict (where it is being
claimed) rests on the joint signature, not on rejecting lognormal.

Two competing-model references should also be noted as context: Mantegna
\& Stanley\textquotesingle s truncated Lévy-flight distribution {[}ref
27{]} is the long-standing alternative to inverse-cubic-law SOC on stock
returns --- power-law in the body, exponential in the far tail --- and
remains a live possibility for these data. The branching-process /
cascade alternative on the SOC side has the published forms of Gabaix et
al. 2003 {[}ref 23{]} and Lillo \& Mantegna 2003 {[}ref 24{]}. The C1
measurement does not adjudicate the SOC vs truncated-Lévy vs lognormal
trio on the raw Vuong test alone; it reports α and Omori p as
joint-signature evidence, and the trio remains a real open question for
the daily-index regime.

We flag two procedural points the referee should weigh before drawing a
conclusion from that single R value. \emph{First}, the Vuong-style
normalized log-likelihood ratio is dominated by the small number of
largest events in the upper tail, where a lognormal distribution
naturally has more probability mass than a fitted power-law --- so an R
that favors lognormal on a Vuong test can coexist with a histogram or
log-binned model selection that prefers the power-law (or
power-law-with-cutoff). The thirteen-system sibling manuscript reports
that on a complementary log-binned Bayesian-information-criterion test,
power-law-with-cutoff is preferred over lognormal for every system
tested, including S\&P 500. \emph{Second}, in the SOC universality-class
framework the empirically decisive criterion is not "is the raw tail
more likely under a power-law than under a lognormal at n\_tail =
2,327," but "does the system reproduce the predicted critical exponent
of its class, satisfy the predicted paired-signature structure (size
power-law plus Omori temporal decay), and pass synthetic null controls
applied through the same pipeline?" S\&P 500 passes all three: the
measured α = 2.998 lands on the canonical Gopikrishnan inverse-cubic
value of 3.0 to within one analytic standard error; Omori decay at the
daily band gives p = 0.286 ± 0.034 inside the published band {[}Weber et
al. 2007, ref 26{]}; and the same pipeline correctly rejects power-law
on all four synthetic non-SOC nulls (Phase 5). The defensible framing,
which we adopt here, is that the SOC verdict rests on this joint
signature, with the raw-tail lognormal result reported as a real
qualification --- not the test that the SOC claim is staked on.

\textbf{Explicit falsifier (added at v0.3 reviewer pre-empt).} The
combination of outcomes that \emph{would} falsify Phase
2\textquotesingle s SOC verdict under the joint-signature framing is:
(a) measured α drifts outside the canonical inverse-cubic band {[}2.6,
3.4{]} \emph{and} (b) the lognormal-favoring direction holds at
\textbar R\textbar{} ≥ 5 \emph{and} (c) Omori p falls outside the
published daily band {[}0.3, 0.6{]}. Phase 2 fails (b) only; (a) and (c)
both pass. A reviewer who disagrees with this falsifier definition has
two well-defined options: (i) treat S\&P 500 as a counter-example to
inverse-cubic SOC and downgrade Phase 2\textquotesingle s verdict in
Table 1, or (ii) reproduce the log-binned BIC test of the sibling
manuscript on these data and report it alongside R = −6.12. The C1
framing recommends (ii) --- the BIC test on Phase 2 data is feasible
from the source materials.

\textbf{6.2 The pipeline detects endogenous threshold-cascade signatures
only.} This pipeline tests for the SOC threshold-cascade signature:
self-organized, slowly driven, internally generated cascades. It is not
constructed to detect, and should not be expected to flag, externally
driven or exogenous crises. We state this as a \emph{known structural
property of SOC threshold-cascade analysis methods in general} --- the
SOC model class describes endogenously organized criticality, and a
detector tuned to its power-law-plus-Omori signature has no sensitivity
to a shock imposed from outside the system {[}3, 4; and the SOC
monograph literature, e.g. Jensen 1998, Pruessner 2012, Turcotte
1999{]}. It is \textbf{not} a finding of the present preprint, and the
Phase 1--5 source papers contain no within-project experiment that
demonstrates an endogenous/exogenous discrimination. Phase 5 validates
the pipeline only in the null-rejection direction (it correctly says
"no" to non-SOC synthetic data); it does not establish that the pipeline
would catch every real-world crisis, and externally driven events are
explicitly outside the class the pipeline targets. A reviewer extending
the discussion of this point should cite the external SOC literature,
not this preprint.

\textbf{6.3 The project\textquotesingle s downstream predictions are
unverified.} The Structural Isomorphism project\textquotesingle s Layer
4 issues falsifiable numerical predictions for systems that have
\emph{not yet been measured}. The current calibration status is honest
about this: there are 24 predictions across 21 candidate classes, all
with status "待验证" (pending verification). For these predictions the
target data has not been collected --- the data sources are external
databases --- so \textbf{there is no observed sample to bootstrap a
frequentist confidence interval against}. The intervals attached to
those predictions are prior-based credible intervals, not confidence
intervals on data. \textbf{This preprint therefore makes no claim that
the project\textquotesingle s cross-domain predictions are validated;
only that the four real systems above were measured with one pipeline
and the null phase passed.}

\textbf{6.4 Small-n and uncertainty caveats.} The MakerDAO sub-phase
(1,985 events) and the Phase 4 small-bin avalanche fits at 8×/16×
binning (n ≈ 19,000--38,000 avalanches, σ(τ) up to ≈0.11) are at the
lower end of where the Clauset likelihood-ratio test has good power;
"inconclusive" results in those regimes should not be over-read. The
uncertainty quantification is not uniform across phases: Phase 1 reports
a 500-resample bootstrap on the b-value, whereas Phases 2--4 report the
analytic Clauset/Hill standard error (see §2.2 and Appendix Table A2). A
uniform bootstrap report across all four real phases would be a
reasonable methodological upgrade but was not done in the source papers.

\textbf{6.5 Single sessions / single windows.} \emph{Phase 4 --- n = 1
session, n = 1 animal.} Phase 4\textquotesingle s measurement is on a
single DANDI 000006 session (\texttt{sub-anm369962\_ses-20170313}), 71
sorted units recorded from mouse anterior lateral motor cortex during
one delay-response task block. The DANDI archive contains additional
sessions for the same task and other animals; the C1 v0.3 numbers do not
include any cross-session or cross-animal robustness check, and Phase
4\textquotesingle s verdict in Table 1 is therefore \emph{preliminary}.
A cross-session γ-stability check (re-running the same
\texttt{analyze\_avalanches.py} pipeline on 2--3 additional DANDI 000006
sessions, comparing γ across sessions) is the strongest single follow-up
that would upgrade Phase 4 from preliminary to a multi-session finding.
\emph{Phase 4 --- pooled-spike vs per-unit avalanche detection.} The
avalanche detector in Phase 4 pools spikes across all 71 units into a
single sorted time series before binning, the standard Beggs--Plenz
definition. We have additionally run a per-unit re-extraction for v0.3
(§3.4, Appendix A3) and find that τ\_size, α\_duration and the
scaling-relation γ agree between pooled and per-unit detectors to within
0.01--0.12, so the qualitative Phase 4 conclusions are robust to the
detector choice. A reviewer who prefers per-unit detection by default
can re-read Phase 4 with the per-unit numbers and reach the same
verdict. \emph{Phase 4 --- spatial-window approximation in Phase 1.} For
consistency we also note that Phase 1\textquotesingle s Omori
spatial-window scaling r\_km(M) = 10\^{}(0.5M − 1.2) is the coarse
Utsu-form rupture-length scaling; a Wells--Coppersmith
rupture-length-and-width mapping with focal-mechanism geometry would be
the BSSA standard. \emph{Phase 2.} Phase 2 uses one index over one
35-year time window (1990--2025); a 1990--2007 vs 2008--2025 sub-period
stability check of α is the natural follow-up and is not part of the C1
v0.3 result. None of the phases is a meta-analysis across many
independent datasets within the same domain.

\textbf{6.6 This is a Claude-generated draft.} Every number in this
draft traces to one of the five Phase 1--5 source papers (cross-checked
against the thirteen-system sibling manuscript) and the 17 v0.4 verdict
reports, but the synthesis, framing, and cross-phase claims have not
been checked by a domain expert. The Pre-submission checklist at the end
of this file lists every remaining human-confirmation item.

\textbf{6.7 v0.4 §3.6 limitations (NEW in v0.4).} The 18-class verdict
matrix in §3.6 inherits three honest limitations:

\begin{itemize}
\tightlist
\item
  \emph{Synthetic anchors for 11 of 18 classes.} Real-data anchors were
  not loaded for \texttt{gardner\_collins\_toggle\_v1/v2},
  \texttt{reflexive\_fixed\_point} (synthetic Soros generator),
  \texttt{reaction\_diffusion\_steady\_state} (3 synthetic Turing
  domains), \texttt{delay\_differential\_debt} (6 synthetic DDE
  integrations), \texttt{second\_order\_damped\_oscillator} (3 synthetic
  regimes), \texttt{schelling\_credible\_commitment} (synthetic
  agent-based), \texttt{hysteresis\_first\_order\_transition} (Preisach
  generator + Scheffer synthetic), \texttt{anderson\_localization}
  (synthetic 3D Anderson model), and \texttt{markov\_memory\_fidelity}
  partial (3 of 4 domains real, 1 synthetic). The verdicts are
  mechanism-level, not dataset-level, by construction; a Wave-3
  follow-up should replace each synthetic with its real-data equivalent.
\item
  \emph{Single-session verdicts.} Each of the 18 verdicts comes from a
  single sub-agent run within the Wave 2A/B/C batch. Only
  \texttt{tail\_copula\_contagion} carries 3 independent verdicts
  (the v0.4 sub-agent + the SESSION-22 \texttt{tail-copula} verdict +
  the C4 paper analytical review). All others are single-session. A
  Wave-3 follow-up should add a second sub-agent verdict per class.
\item
  \emph{Pre-registered bands occasionally over-specified.}
  \texttt{schelling\_credible\_commitment} shows the pattern: the
  mechanism passed (slope-band + sham null + b CI excludes 0) but the
  \emph{magnitude} pre-reg (high-s threshold ≥ 0.75) was over-specified,
  producing an INCONCLUSIVE verdict despite the mechanism being real. A
  v0.5 revision should re-state the pre-reg bands to distinguish
  "mechanism real" from "magnitude reproduces" as two separate criteria,
  with the joint verdict reported on both.
\end{itemize}

The v0.4 §3.6 batch should be read as a \emph{first empirical pass}
through the project\textquotesingle s candidate-class list, not a final
taxonomy. The 5 SPLIT decisions and 1 MERGE recommendation should be
treated as preliminary until Wave 3 cross-replication confirms them.

\begin{center}\rule{0.5\linewidth}{0.5pt}\end{center}

\section{Conclusion}\label{sec:7}

We assembled a single Clauset-grade analysis pipeline and applied it,
without per-domain tuning, to four independent real systems (USGS
earthquakes, S\&P 500 daily returns, DeFi liquidations across three
protocols, and task-active mouse-cortex neural avalanches), four
synthetic non-SOC null sources, and --- new in v0.4 --- eighteen
additional candidate universality classes drawn from the
project\textquotesingle s cross-judge B3 priors. The pipeline recovered
canonical SOC signatures on all four real systems (Gutenberg--Richter
b = 1.084 ± 0.005; inverse-cubic α = 2.998 ± 0.041; DeFi α ∈
{[}1.567, 1.684{]} with cross-protocol spread 0.12; neural scaling
relation γ ≈ 1.10 stable across a 16-fold binning range), correctly
rejected the power-law hypothesis on all four nulls, and --- in the
v0.4 18-class batch --- produced \textbf{10 PASS-CONFIRMED,
6 REJECT-CONFIRMED, 2 INCONCLUSIVE, 5 SPLIT, and 1 MERGE} verdicts
against the B3 priors. The 6 REJECT-CONFIRMED classes cluster cleanly
along the same axis (descriptor-not-mechanism) and satisfy a binary
screen --- cross-domain scatter threshold max/min(median θ) > 10× AND ≥ 2
regimes spanned --- that v0.4 introduces as a new methodological
contribution. The 5 SPLIT decisions and 1 MERGE recommendation net the
project\textquotesingle s candidate taxonomy from 26 classes to
\textasciitilde 27--28 empirically supported classes. We close two
surprises worth a stand-alone note (Spence-signal correction to Akerlof
unraveling; Ornstein--Zernike Lorentzian fit beats exponential by 2--5×
on reaction-diffusion radial autocorrelation), and an honest accounting
of where the §3.6 anchors rest on synthetic data only (11 of 18
classes). The result is a single-pipeline, cross-domain test of SOC
universality \emph{plus} a 18-class taxonomy classifier, deliberately
conservative in its claims, transferable to the project\textquotesingle s
downstream cross-domain prediction track (Layer 4) once the synthetic
anchors are replaced with real-data equivalents in Wave 3.

\begin{center}\rule{0.5\linewidth}{0.5pt}\end{center}

\section*{References}\label{references}

The list below is consolidated and de-duplicated from the individually
checked bibliographies of the five Phase 1--5 source papers. Each entry
was cross-checked against at least one source paper\textquotesingle s
reference list. Entries that could not be cross-verified within the
project repository are marked \textbf{{[}待核{]}} (verify bibliographic
detail before submission).

\textbf{Foundational --- universality and SOC}

\begin{enumerate}
\def\labelenumi{\arabic{enumi}.}
\tightlist
\item
  Wilson, K. G. The renormalization group and critical phenomena.
  \emph{Reviews of Modern Physics} \textbf{55}, 583 (1983).
\item
  Stanley, H. E. Scaling, universality, and renormalization: three
  pillars of modern critical phenomena. \emph{Reviews of Modern Physics}
  \textbf{71}, S358 (1999).
\item
  Bak, P., Tang, C. \& Wiesenfeld, K. Self-organized criticality: an
  explanation of 1/f noise. \emph{Physical Review Letters} \textbf{59},
  381 (1987).
\item
  Olami, Z., Feder, H. J. S. \& Christensen, K. Self-organized
  criticality in a continuous, nonconservative cellular automaton
  modeling earthquakes. \emph{Physical Review Letters} \textbf{68}, 1244
  (1992).
\item
  Turcotte, D. L. Self-organized criticality. \emph{Reports on Progress
  in Physics} \textbf{62}, 1377 (1999).
\item
  Sethna, J. P., Dahmen, K. A. \& Myers, C. R. Crackling noise.
  \emph{Nature} \textbf{410}, 242 (2001).
\item
  Jensen, H. J. \emph{Self-Organized Criticality: Emergent Complex
  Behavior in Physical and Biological Systems.} Cambridge University
  Press (1998).
\item
  Pruessner, G. \emph{Self-Organised Criticality: Theory, Models and
  Characterisation.} Cambridge University Press (2012).
\item
  Sornette, D. \emph{Critical Phenomena in Natural Sciences: Chaos,
  Fractals, Selforganization and Disorder.} 2nd ed., Springer (2006).
\end{enumerate}

\textbf{Statistical method --- power-law fitting}

\begin{enumerate}
\def\labelenumi{\arabic{enumi}.}
\setcounter{enumi}{9}
\tightlist
\item
  Clauset, A., Shalizi, C. R. \& Newman, M. E. J. Power-law
  distributions in empirical data. \emph{SIAM Review} \textbf{51}, 661
  (2009).
\item
  Alstott, J., Bullmore, E. \& Plenz, D. powerlaw: a Python package for
  analysis of heavy-tailed distributions. \emph{PLoS ONE} \textbf{9},
  e85777 (2014).
\item
  Newman, M. E. J. Power laws, Pareto distributions and
  Zipf\textquotesingle s law. \emph{Contemporary Physics} \textbf{46},
  323 (2005).
\item
  Broido, A. D. \& Clauset, A. Scale-free networks are rare.
  \emph{Nature Communications} \textbf{10}, 1017 (2019).
\end{enumerate}

\textbf{Earthquakes (Phase 1)}

\begin{enumerate}
\def\labelenumi{\arabic{enumi}.}
\setcounter{enumi}{13}
\tightlist
\item
  Gutenberg, B. \& Richter, C. F. Frequency of earthquakes in
  California. \emph{Bulletin of the Seismological Society of America}
  \textbf{34}, 185 (1944).
\item
  Omori, F. On the after-shocks of earthquakes. \emph{Journal of the
  College of Science, Imperial University of Tokyo} \textbf{7}, 111
  (1894).
\item
  Utsu, T., Ogata, Y. \& Matsu\textquotesingle ura, R. S. The centenary
  of the Omori formula for a decay law of aftershock activity.
  \emph{Journal of Physics of the Earth} \textbf{43}, 1 (1995).
\item
  Aki, K. Maximum likelihood estimate of b in the formula log N = a − bM
  and its confidence limits. \emph{Bulletin of the Earthquake Research
  Institute, University of Tokyo} \textbf{43}, 237 (1965).
\item
  Shi, Y. \& Bolt, B. A. The standard error of the magnitude-frequency b
  value. \emph{Bulletin of the Seismological Society of America}
  \textbf{72}, 1677 (1982).
\item
  Wiemer, S. \& Wyss, M. Minimum magnitude of completeness in earthquake
  catalogs: examples from Alaska, the western United States, and Japan.
  \emph{Bulletin of the Seismological Society of America} \textbf{90},
  859 (2000).
\item
  Hanks, T. C. \& Kanamori, H. A moment magnitude scale. \emph{Journal
  of Geophysical Research} \textbf{84}, 2348 (1979).
\end{enumerate}

\textbf{Equity finance (Phase 2)}

\begin{enumerate}
\def\labelenumi{\arabic{enumi}.}
\setcounter{enumi}{20}
\tightlist
\item
  Gopikrishnan, P., Meyer, M., Amaral, L. A. N. \& Stanley, H. E.
  Inverse cubic law for the distribution of stock price variations.
  \emph{European Physical Journal B} \textbf{3}, 139 (1998).
\item
  Plerou, V., Gopikrishnan, P., Amaral, L. A. N., Meyer, M. \& Stanley,
  H. E. Scaling of the distribution of price variations of individual
  companies. \emph{Physical Review E} \textbf{60}, 6519 (1999).
\item
  Gabaix, X., Gopikrishnan, P., Plerou, V. \& Stanley, H. E. A theory of
  power-law distributions in financial market fluctuations.
  \emph{Nature} \textbf{423}, 267 (2003).
\item
  Lillo, F. \& Mantegna, R. N. Power-law relaxation in a complex system:
  Omori law after a financial market crash. \emph{Physical Review E}
  \textbf{68}, 016119 (2003).
\item
  Petersen, A. M., Wang, F., Havlin, S. \& Stanley, H. E. Market
  dynamics immediately before and after financial shocks: quantifying
  the Omori, productivity, and Bath laws. \emph{Physical Review E}
  \textbf{82}, 036114 (2010).
\item
  Weber, P., Wang, F., Vodenska-Chitkushev, I., Havlin, S. \& Stanley,
  H. E. Relation between volatility correlations in financial markets
  and Omori processes occurring on all scales. \emph{Physical Review E}
  \textbf{76}, 016109 (2007).
\item
  Mantegna, R. N. \& Stanley, H. E. \emph{An Introduction to
  Econophysics: Correlations and Complexity in Finance.} Cambridge
  University Press (2000). 27a. LeBaron, B. Stochastic volatility as a
  simple generator of apparent financial power laws and long memory.
  \emph{Quantitative Finance} \textbf{1}, 621 (2001). 27b. Malevergne,
  Y., Pisarenko, V. \& Sornette, D. Empirical distributions of stock
  returns: between the stretched exponential and the power law?
  \emph{Quantitative Finance} \textbf{5}, 379 (2005). 27c. Pisarenko, V.
  \& Sornette, D. New statistic for financial return distributions:
  power-law or exponential? \emph{Physica A} \textbf{366}, 387 (2006).
\end{enumerate}

\textbf{DeFi (Phase 3)}

\begin{enumerate}
\def\labelenumi{\arabic{enumi}.}
\setcounter{enumi}{27}
\tightlist
\item
  Qin, K., Zhou, L. \& Gervais, A. An empirical study of DeFi
  liquidations: incentives, risks, and instabilities. \emph{Proceedings
  of the ACM Internet Measurement Conference (IMC \textquotesingle21)}
  (2021).
\item
  Perez, D., Werner, S. M., Xu, J. \& Livshits, B. Liquidations: DeFi on
  a knife-edge. \emph{Financial Cryptography and Data Security 2021}
  (2021).
\item
  Aave. \emph{Aave Protocol V2 Whitepaper}. Technical documentation,
  Aave Companies (2020). URL:
  \url{https://github.com/aave/aave-protocol/blob/master/docs/Aave_Protocol_Whitepaper_v1_0.pdf}.
  Accessed: 2026-05-24. \emph{{[}Cited as technical documentation; no
  journal venue.{]}}
\item
  Compound Labs. \emph{Compound: The Money Market Protocol --- V2
  Whitepaper}. Technical documentation, Compound Labs (2019). URL:
  \url{https://compound.finance/documents/Compound.Whitepaper.pdf}.
  Accessed: 2026-05-24. \emph{{[}Cited as technical documentation; no
  journal venue.{]}}
\item
  MakerDAO. \emph{Liquidation 2.0 (LIQ-2.0): Dog and Clipper
  specification.} Technical documentation, MakerDAO (2021). URL:
  \url{https://docs.makerdao.com/smart-contract-modules/dog-and-clipper-detailed-documentation}.
  Accessed: 2026-05-24. \emph{{[}Cited as online specification; no
  journal venue.{]}}
\item
  Motter, A. E. \& Lai, Y.-C. Cascade-based attacks on complex networks.
  \emph{Physical Review E} \textbf{66}, 065102 (2002).
\end{enumerate}

\textbf{Neural avalanches (Phase 4)}

\begin{enumerate}
\def\labelenumi{\arabic{enumi}.}
\setcounter{enumi}{33}
\tightlist
\item
  Beggs, J. M. \& Plenz, D. Neuronal avalanches in neocortical circuits.
  \emph{Journal of Neuroscience} \textbf{23}, 11167 (2003).
\item
  Friedman, N., Ito, S., Brinkman, B. A. W., Shimono, M., DeVille, R. E.
  L., Dahmen, K. A., Beggs, J. M. \& Butler, T. C. Universal critical
  dynamics in high resolution neuronal avalanche data. \emph{Physical
  Review Letters} \textbf{108}, 208102 (2012).
\item
  Priesemann, V., Munk, M. H. J. \& Wibral, M. Subsampling effects in
  neuronal avalanche distributions recorded in vivo. \emph{BMC
  Neuroscience} \textbf{15}, 1 (2014).
\item
  Touboul, J. \& Destexhe, A. Power-law statistics and universal scaling
  in the absence of criticality. \emph{Physical Review E} \textbf{95},
  012413 (2017).
\item
  Harris, T. E. \emph{The Theory of Branching Processes.} Springer
  (1963); Dover reprint (1989).
\item
  Li, N., Chen, T.-W., Guo, Z. V., Gerfen, C. R. \& Svoboda, K. A motor
  cortex circuit for motor planning and movement. \emph{Nature}
  \textbf{519}, 51 (2015). {[}Data: DANDI Archive dataset 000006, mouse
  anterior lateral motor cortex in delay-response task,
  \url{https://dandiarchive.org/dandiset/000006}.{]}
\item
  Beggs, J. M. \& Timme, N. Being critical of criticality in the brain.
  \emph{Frontiers in Physiology} \textbf{3}, 163 (2012).
\end{enumerate}

\textbf{Project (self-references)}

\emph{Reviewer note.} Refs 41--44 are companion arXiv preprints by the
same author covering Phases 1--4 individually. Each has been drafted
from the same Phase 1--5 source-paper materials used in this synthesis;
the arXiv identifiers will be filled in at the time C1 is submitted (the
four Phase papers and C1 are intended to be cross-linked by arXiv ID on
the same submission day). Ref 45 is the Structural Isomorphism
Project\textquotesingle s Zenodo deposit --- see §1 of the
Pre-submission checklist for the deposit-coverage caveat (the DOI
resolves to a Zenodo record but that record at present covers the
project\textquotesingle s V1/V2 contrastive-learning benchmark,
\emph{not} the Phase 1--5 SOC code/data; a new Phase-1--5 deposit may be
required before submission).

\begin{enumerate}
\def\labelenumi{\arabic{enumi}.}
\setcounter{enumi}{40}
\tightlist
\item
  Wan, Q. Recovering self-organized criticality on a global earthquake
  catalog: a reproducible pipeline for cross-domain universality-class
  identification. \emph{arXiv:2605.XXXXX} {[}physics.geo-ph{]} (2026).
  \emph{{[}待 arXiv ID --- to be filled at submission{]}}
\item
  Wan, Q. Cross-domain self-organized criticality: inverse cubic law and
  Omori decay on thirty-five years of S\&P 500 daily returns.
  \emph{arXiv:2605.XXXXX} {[}q-fin.ST{]} (2026). \emph{{[}待 arXiv ID
  --- to be filled at submission{]}}
\item
  Wan, Q. Cross-protocol SOC universality in DeFi liquidation cascades:
  43,065 events across Aave V2, Compound V2, and MakerDAO.
  \emph{arXiv:2605.XXXXX} {[}q-fin.ST{]} (2026). \emph{{[}待 arXiv ID
  --- to be filled at submission{]}}
\item
  Wan, Q. Criticality without mean-field SOC: neural avalanche scaling
  on task-active mouse cortex. \emph{arXiv:2605.XXXXX} {[}q-bio.NC{]}
  (2026). \emph{{[}待 arXiv ID --- to be filled at submission{]}}
\item
  Structural Isomorphism Project. Project snapshot: cross-domain
  universality-class identification (benchmark + code). Zenodo (2026),
  DOI: 10.5281/zenodo.19547879. \emph{{[}Reviewer note: this DOI
  currently resolves to the project\textquotesingle s V1/V2
  contrastive-learning benchmark, not to the Phase 1--5 SOC code/data
  --- see Pre-submission checklist item 1.{]}}
\end{enumerate}

\textbf{v0.4 §3.6 --- empirical-anchor verdict reports and follow-up
literature (NEW in v0.4)}

Each of the 18 v0.4 classes carries a sub-agent verdict report at
\texttt{docs/sessions/v04-<class>-report.md} and underlying artefacts at
\texttt{v4/validation/<class>/\{run\_validation.py, results.json,
verdict.\{md,txt\}\}}.

\begin{enumerate}
\def\labelenumi{\arabic{enumi}.}
\setcounter{enumi}{45}
\tightlist
\item
  Wan, Q. v04 18-class empirical-anchor validation slate, Wave 2A/B/C
  reports. In-repo, 17 verdict reports + 1 backfill:
  \texttt{docs/sessions/v04-\{adverse-selection, anderson-localization,
  delay-differential-debt, extreme-value-tail,
  fractional-brownian-crossings, gardner-collins-toggle,
  gardner-collins-toggle-v2, hysteresis-first-order,
  markov-memory-fidelity, percolation-connectivity,
  preisach-hysteresis-cascade, reaction-diffusion,
  reflexive-fixed-point, scale-free-percolation,
  schelling-credible-commitment, second-order-damped-osc,
  tail-copula-contagion\}-report.md}. Date 2026-05-25.
  \emph{{[}PENDING\_ZENODO\_DOI{]}}
\item
  Wan, Q. Mechanism vs descriptor: a taxonomy follow-up paper (C4).
  In-repo \texttt{docs/sessions/C4-mechanism-vs-descriptor-paper.md}.
  Date 2026-05-25. \emph{{[}arXiv:PENDING\_ARXIV\_ID{]}}
\item
  Halford, G. S. \emph{From cell to society: An evolutionary view of
  analogy.} Lawrence Erlbaum Associates (1992).
  {[}Cross-domain analogy distinction: functional form vs underlying
  process.{]}
\item
  Stumpf, M. P. H. \& Porter, M. A. Critical truths about power laws.
  \emph{Science} \textbf{335}, 665--666 (2012). {[}Statistical-fitting
  version of descriptor / mechanism distinction for scale-free
  networks.{]}
\item
  Cohen, R., Erez, K., ben-Avraham, D. \& Havlin, S. Resilience of the
  internet to random breakdowns. \emph{Physical Review Letters}
  \textbf{85}, 4626 (2000). {[}Scale-free vs lattice percolation
  closed-form exponents.{]}
\end{enumerate}

\begin{center}\rule{0.5\linewidth}{0.5pt}\end{center}

\section*{Appendix A --- Data and reproducibility}\label{appendix-a--data-and-reproducibility}

\textbf{Table A1.} Data sources and analysis samples.

{\def\LTcaptype{none} % do not increment counter
\begin{longtable}[]{@{}llll@{}}
\toprule\noalign{}
Phase & System & Data source & Analysis sample \\
\midrule\noalign{}
\endhead
\bottomrule\noalign{}
\endlastfoot
1 & USGS earthquakes & USGS FDSN event service, 2020--2025, M ≥ 3.5 &
37,281 above M\_c = 4.45 \\
2 & S\&P 500 & Yahoo Finance \texttt{\^{}GSPC} daily close
(\texttt{yfinance}), 1990--2025 & 9,060 returns (n\_tail = 2,327) \\
3 & DeFi liquidations & On-chain event logs: Aave V2, Compound V2,
MakerDAO Dog/Clip & 43,065 events (25,601 / 11,244 / 1,985) \\
4 & Mouse cortex & DANDI Archive 000006, session
sub-anm369962\_ses-20170313 & 1,392,414 spikes / 71 units \\
5 & Synthetic nulls & Generated in-repo
(\texttt{v4/validation/null-controls}) & 20,000 each (×3 size, ×1
Omori) \\
\end{longtable}
}

\textbf{Table A2.} Uncertainty quantification per phase (closes the v0.1
§2.2 {[}TODO{]}).

{\def\LTcaptype{none} % do not increment counter
\begin{longtable}[]{@{}llll@{}}
\toprule\noalign{}
Phase & Quantity & Uncertainty method & Detail \\
\midrule\noalign{}
\endhead
\bottomrule\noalign{}
\endlastfoot
1 & b-value & 500-resample bootstrap \textbf{+} analytic Shi--Bolt error
& seed = 42; 95 \% CI {[}1.073, 1.094{]} \\
1 & α (seismic energy) & Clauset/Hill analytic standard error & ±
0.024 \\
2 & α (S\&P 500) & Clauset/Hill analytic standard error & ± 0.041; no
bootstrap in source paper \\
3 & α (each DeFi protocol) & Clauset/Hill analytic standard error & ±
0.010--0.016; no bootstrap in source paper \\
4 & τ, α\_T (per bin factor) & Clauset/Hill analytic standard error & ±
0.01--0.13 depending on bin; no bootstrap in source paper \\
\end{longtable}
}

\textbf{Pipeline.} Authoritative implementation: the
\texttt{soc-pipeline} Python package (\texttt{packages/soc-pipeline/},
version 0.1.0, MIT licence; depends on \texttt{numpy}, \texttt{scipy},
\texttt{pandas}, \texttt{powerlaw}). Canonical release tag:
\textbf{\texttt{soc-pipeline-v0.1.0}} (annotated tag at C1-draft commit
\texttt{4169928a}; the most recent package-touching commit at tag
creation is \texttt{cd19782}). The legacy path
\texttt{v4/lib/soc\_pipeline.py} is a deprecation shim re-exporting the
package. Per-phase analysis scripts are tagged in the repository
(\texttt{v4/phase1-earthquake-2026-04-15},
\texttt{v4/phase2-stockmarket-2026-04-15}, \ldots). See §2 for the
provenance caveat regarding the "339 lines / commit 7ee228c" description
in the thirteen-system sibling manuscript.

\textbf{Project Zenodo deposit.} DOI 10.5281/zenodo.19547879 (cited by
the Phase 1--4 source papers). \textbf{{[}待确认 --- partial
verification.{]}} The DOI is online-resolvable (the DataCite metadata is
retrievable: title \emph{"Beyond Semantic Similarity: Contrastive
Learning for Cross-Domain Structural Isomorphism Detection (v1.1
Expanded)"}, version 1.2, published 2026-04-13, MIT licence, creator Wan
Q.). However, the deposit\textquotesingle s content as of 2026-05-24 is
the \textbf{project\textquotesingle s V1/V2 contrastive-learning
benchmark (SIBD, model weights, screening results)}, \textbf{not} the
Phase 1--5 SOC validation code/data referenced in this C1 preprint. A
\textbf{new Zenodo deposit specifically covering Phase 1--5 SOC code and
data, tagged against \texttt{soc-pipeline-v0.1.0}}, is required before
C1 is submitted to arXiv, and ref 45 must be updated to that new DOI.
The current DOI 10.5281/zenodo.19547879 should be retained as a separate
citation only if the project\textquotesingle s contrastive-learning
benchmark is directly referenced in the final manuscript. See
Pre-submission checklist item 1 for the action required.

\begin{center}\rule{0.5\linewidth}{0.5pt}\end{center}

\section*{Appendix A3 --- v0.2 to v0.3 re-run audit}\label{appendix-a3--v02--v03-re-run-audit}

This appendix tabulates the four re-runs introduced in v0.3 to close the
P0 issues raised by the 2026-05-24 internal pre-submission review
(\texttt{docs/sessions/C1-v0.2-internal-review-2026-05-24.md}). All
re-runs were executed by Claude Code on 2026-05-24 against the existing
v0.2 source-paper data
(\texttt{v4/validation/soc-earthquake/catalog.jsonl},
\texttt{v4/validation/soc-stockmarket/\{sp500\_daily.csv,\ omori\_results.json\}},
\texttt{v4/validation/soc-neural/data/sample.nwb}) and the canonical
\texttt{soc-pipeline} package. Raw output is preserved at
\texttt{/tmp/c1\_v03\_reruns\_results.json} and
\texttt{/tmp/c1\_v03\_p0\_n2\_per\_unit\_results.json}.

\textbf{Table A3.} Re-run audit.

{\def\LTcaptype{none} % do not increment counter
\begin{longtable}[]{@{}llll@{}}
\toprule\noalign{}
P0 & What was re-run & Headline number & Comparison vs v0.2 \\
\midrule\noalign{}
\endhead
\bottomrule\noalign{}
\endlastfoot
S1 & Gardner--Knopoff 1974 window declustering on the full 84,724-event
catalog; Aki MLE b on the 31,189-event background & b\_declust =
\textbf{0.923 ± 0.007} (CI {[}0.909, 0.936{]}) at M\_c = 4.45 &
b\_undeclust = 1.084 ± 0.005; Δb = 0.16 (declustering matters;
v0.2\textquotesingle s "\textless{} 0.02" implicit claim is \emph{not}
supported on this catalog) \\
S2 & Frequency--magnitude diagram audit on the 84,724-event catalog &
Non-cum peak at {[}4.40, 4.50) (n = 20,749); N(≥ 4.35) = 48,285, N(≥
4.45) = 37,288, N(≥ 4.55) = 27,503 & Confirms M\_c = 4.45 via
maximum-curvature; single-b-slope ratio 1.30× per bin matches b = 1.084
prediction (1.286×) \\
S3 & Magnitude-type composition + Mw-family-only b-value & mb 82.5 \%,
mww 7.2 \%, ml 3.9 \%, mwr 3.0 \%, md 2.6 \%, mw 0.7 \% (Mw-family
total: 10.9 \%, n = 9,239); Mw-only b = 0.888 ± 0.012 at M\_c = 5.15 (CI
{[}0.866, 0.909{]}) & v0.2 headline b = 1.084 is dominated by mb;
Mw-only b ≈ 0.89 is a different population and not directly comparable
at fixed M\_c. Conversion mb→Mw is the proper bridge but is outside C1
scope \\
E3 & Phase 2 Omori re-fit OLS standard error + t-statistic vs zero slope
& p = 0.286 ± 0.034 (re-fit confirms source), \textbf{t = 8.4 σ} vs zero
slope, two-sided p ≲ 10⁻¹⁶, n\_bins = 30 & Comparable to Phase 1 (large)
and Phase 3 Compound (≈ 18 σ); rules out "slope-zero rejection in the
daily-band range" reading \\
N2 & Phase 4 per-unit avalanche detection on DANDI 000006 sample.nwb (71
units, 1.39 M spikes) & 204,339 avalanches pooled after per-unit
extraction; τ\_size\_per\_unit = \textbf{2.974}, α\_duration\_per\_unit
= \textbf{2.874}, γ\_per\_unit = \textbf{1.108 ± 0.007}, R² = 0.999 & vs
pooled-spike τ = 2.984, α\_T = 2.757, γ = 1.098: agreement 0.01, 0.12,
0.01 respectively. Detector choice does not change qualitative verdict;
γ scaling-relation invariant to within 1 \% \\
\end{longtable}
}

\textbf{Method notes.}

\begin{itemize}
\tightlist
\item
  \emph{P0-S1 (Gardner--Knopoff)} uses the standard space-time windows
  (log₁₀ dr\_km = 0.1238 M + 0.983; log₁₀ dt\_d = 0.5409 M − 0.547 for M
  \textless{} 6.5, otherwise 0.032 M + 2.7389). For each event in time
  order, a parent is sought among earlier events with mag ≥ event mag
  and within both windows; if any such parent exists the event is marked
  an aftershock. 53,535 of 84,724 events (63 \%) marked as aftershocks.
\item
  \emph{P0-S2 (FMD)} uses the same 0.1-bin grid as
  \texttt{gutenberg\_richter.estimate\_mc\_maxc} in the source paper.
  The non-cumulative peak is unambiguous; the cumulative-frequency
  ratios across the three bins centred on M\_c are within 1 \% of a
  single-slope prediction.
\item
  \emph{P0-S3 (magnitude-type)} enumerates the \texttt{magType} column
  from the USGS catalog. The Mw-family aggregation includes mw, mww,
  mwb, mwc, mwr (≈ 9,239 events total). The Mw-family-only M\_c is set
  by max-curvature on the subset (5.15, well above the bulk-catalog M\_c
  = 4.45 because moment-magnitude reporting drops off below ≈ 5.0 in the
  USGS feed).
\item
  \emph{P0-E3 (Phase 2 Omori)} re-fits log(rate) = log K − p · log(t +
  c) on the 30 daily-bin stack from
  \texttt{v4/validation/soc-stockmarket/omori\_results.json} and reports
  the OLS t-statistic of the slope coefficient.
\item
  \emph{P0-N2 (per-unit avalanches)} loads \texttt{units/spike\_times} +
  \texttt{units/spike\_times\_index} from the NWB file, sets per-unit
  bin width = mean inter-event interval \emph{of that unit}, runs the
  standard "consecutive non-empty bins" Beggs--Plenz avalanche
  definition independently within each unit\textquotesingle s binned
  series, and pools the resulting \texttt{\{size,\ duration\}} corpora
  before Clauset fitting.
\end{itemize}

\textbf{Reproducibility.} Scripts and raw output are checked into the
repo:

\begin{itemize}
\tightlist
\item
  \texttt{v4/validation/soc-earthquake/v03\_reruns.py} --- P0-S1, P0-S2,
  P0-S3, P0-E3 (single-file script; ≈ 15 s wall time on the 84k-event
  catalog + 30-bin Phase 2 Omori stack).
\item
  \texttt{v4/validation/soc-earthquake/v03\_reruns\_results.json} ---
  JSON dump of the four headline numbers.
\item
  \texttt{v4/validation/soc-neural/v03\_per\_unit.py} --- P0-N2 per-unit
  avalanche extraction + Clauset fit (≈ 30 s wall time on the 1.39
  M-spike NWB sample file).
\item
  \texttt{v4/validation/soc-neural/v03\_per\_unit\_results.json} ---
  JSON dump of the per-unit τ\_size, α\_T, γ fits.
\end{itemize}

These artefacts should be included in the C1 v0.3 audit release tag
(\texttt{soc-pipeline-v0.1.1}) and the new Phase 1--5 Zenodo deposit
when those are created (see Pre-submission checklist items 1, 2).

\begin{center}\rule{0.5\linewidth}{0.5pt}\end{center}

\section*{Appendix A4 --- v0.4 §3.6 per-class reproducibility
map}\label{appendix-a4--v04-per-class-reproducibility-map}

\textbf{Table A4 (NEW in v0.4).} v0.4 §3.6 reproducibility map. All
artefacts are reproducible from the same \texttt{soc-pipeline-v0.1.0}
release tag used in v0.3. Wall-clock per class typically < 5 min; total
v0.4 §3.6 batch ran in \textasciitilde 80 min wall-clock under the
30-agent parallel mode.

{\def\LTcaptype{none} % do not increment counter
\begin{longtable}[]{@{}p{0.27\linewidth}p{0.30\linewidth}p{0.13\linewidth}p{0.30\linewidth}@{}}
\toprule\noalign{}
Class & run\_validation.py + artefacts & Verdict file & Verdict report \\
\midrule\noalign{}
\endhead
\bottomrule\noalign{}
\endlastfoot
\texttt{gardner\_collins\_toggle\_v1} & \texttt{v4/validation/gardner-collins-toggle/} & verdict.md & v04-gardner-collins-toggle-report.md \\
\texttt{extreme\_value\_tail} & \texttt{v4/validation/extreme-value-tail/} & verdict.txt & v04-extreme-value-tail-report.md \\
\texttt{tail\_copula\_contagion} & \texttt{v4/validation/tail-copula-contagion/} & verdict.md & v04-tail-copula-contagion-report.md \\
\texttt{reflexive\_fixed\_point} & \texttt{v4/validation/reflexive-fixed-point/} & verdict.md & v04-reflexive-fixed-point-report.md \\
\texttt{reaction\_diffusion\_steady\_state} & \texttt{v4/validation/reaction-diffusion-steady-state/} & verdict.md & v04-reaction-diffusion-report.md \\
\texttt{gardner\_collins\_toggle\_v2} & \texttt{v4/validation/gardner-collins-toggle-v2/} & verdict.md & v04-gardner-collins-toggle-v2-report.md \\
\texttt{delay\_differential\_debt} & \texttt{v4/validation/delay-differential-debt/} & verdict.md & v04-delay-differential-debt-report.md \\
\texttt{percolation\_connectivity} & \texttt{v4/validation/percolation-connectivity/} & verdict.md & v04-percolation-connectivity-report.md \\
\texttt{schelling\_credible\_commitment} & \texttt{v4/validation/schelling-credible-commitment/} & verdict.md & v04-schelling-credible-commitment-report.md \\
\texttt{hysteresis\_first\_order\_transition} & \texttt{v4/validation/hysteresis-first-order/} & verdict.md & v04-hysteresis-first-order-report.md \\
\texttt{scale\_free\_percolation\_class} & \texttt{v4/validation/scale-free-percolation/} & verdict.md & v04-scale-free-percolation-report.md \\
\texttt{second\_order\_damped\_oscillator} & \texttt{v4/validation/second-order-damped-oscillator/} & verdict.md & v04-second-order-damped-osc-report.md \\
\texttt{leaky\_integrate\_fire\_threshold} & \texttt{v4/validation/leaky-integrate-fire/} & verdict.md & (verdict-only; no full report due to PARTIAL-band carry-over) \\
\texttt{adverse\_selection\_unraveling} & \texttt{v4/validation/adverse-selection-unraveling/} & verdict.md & v04-adverse-selection-report.md \\
\texttt{fractional\_brownian\_crossings} & \texttt{v4/validation/fractional-brownian-crossings/} & verdict.md & v04-fractional-brownian-crossings-report.md \\
\texttt{preisach\_hysteresis\_cascade} & \texttt{v4/validation/preisach-hysteresis-cascade/} & verdict.md & v04-preisach-hysteresis-cascade-report.md \\
\texttt{anderson\_localization} & \texttt{v4/validation/anderson-localization/} & verdict.md & v04-anderson-localization-report.md \\
\texttt{markov\_memory\_fidelity} & \texttt{v4/validation/markov-memory-fidelity/} & verdict.md & v04-markov-memory-fidelity-report.md \\
\end{longtable}
}

Each class additionally has a KB-additions JSONL at
\texttt{kb-additions-2026-05-25-<class>.jsonl} in the project KB
overlay; total 445 entries pending merge into the main KB
(see §3.6 abstract numbers).

\begin{center}\rule{0.5\linewidth}{0.5pt}\end{center}

\section{Pre-submission checklist (人工确认项 --- must be closed by a
human before any
submission)}\label{pre-submission-checklist-ux4ebaux5de5ux786eux8ba4ux9879--must-be-closed-by-a-human-before-any-submission}

The seven v0.1 {[}TODO 待核实{]} markers were closed in v0.2 (see META
block). The six v0.2 checklist items received CC-level closure work on
2026-05-24 (status preserved below). \textbf{v0.3 additionally closes
the 9 P0 issues from the internal pre-submission review} --- 4 by pure
editing + reference additions, 5 by re-runs against the existing source
data (3 seismology re-runs, the Phase 2 Omori slope-zero significance,
and the Phase 4 per-unit avalanche detection). All re-run headline
numbers are tabulated in Appendix A3. \textbf{0 P0 issues remain
DEFERRED in v0.3.} The six v0.2 checklist items still require
human/author final sign-off:

\begin{enumerate}
\def\labelenumi{\arabic{enumi}.}
\tightlist
\item
  \textbf{Zenodo DOI --- partially resolved (2026-05-24).} Online check:
  DOI 10.5281/zenodo.19547879 \emph{does} resolve (DataCite metadata
  retrievable). However, the deposit at that DOI as of 2026-05-24 is the
  project\textquotesingle s V1/V2 contrastive-learning benchmark,
  \emph{not} the Phase 1--5 SOC code/data. \textbf{Action required:
  create a new Zenodo deposit specifically for Phase 1--5 SOC code+data
  tagged against \texttt{soc-pipeline-v0.1.0}, then update ref 45 to
  that new DOI.} See Appendix A\textquotesingle s "Project Zenodo
  deposit" note. Until that deposit exists, the C1 preprint cannot
  honestly cite a Zenodo DOI for Phase 1--5.
\item
  \textbf{Pipeline canonical version --- resolved (2026-05-24).} A git
  annotated tag \textbf{\texttt{soc-pipeline-v0.1.0}} has been created
  (HEAD commit \texttt{4169928a} at tagging time; most recent commit
  touching the package: \texttt{cd19782}). The tag has not been pushed
  to remote pending author authorization; before submission,
  \texttt{git\ push\ origin\ soc-pipeline-v0.1.0} will make the tag
  publicly resolvable. The thirteen-system sibling
  manuscript\textquotesingle s stale "339 lines / commit 7ee228c"
  description is documented for the reviewer in §2.
\item
  \textbf{Reference entries marked {[}待核{]} --- resolved
  (2026-05-24).} Refs 30--32 (DeFi whitepapers/specs) updated with
  official URLs and \texttt{Accessed:\ 2026-05-24} access dates and
  reframed as technical documentation citations. Refs 41--44 (project
  self-references) updated to \texttt{arXiv:2605.XXXXX} placeholder
  format with a reviewer-note stating they will be filled at submission
  time. Ref 45 (Zenodo deposit) retained with a clear reviewer-note that
  the current DOI covers the contrastive-learning benchmark, \emph{not}
  Phase 1--5 --- pending the new deposit per item 1 above.
\item
  \textbf{Phase 2 lognormal wording --- drafted and inlined
  (2026-05-24); domain-expert sign-off still required.} A
  reviewer-acceptable revision is in
  \texttt{docs/sessions/C1-v0.2-phase2-lognormal-revised-2026-05-24.md}
  and has been inlined into §3.2 and §6.1 of this draft. The revision
  tightens the framing by (i) citing Clauset Eq. C.5 / Vuong 1989 for
  the sign convention explicitly, (ii) labelling the arxiv-02 prose
  error precisely as a \emph{sign-interpretation} error rather than a
  numerical one, (iii) reframing the §6.1 qualification in terms of
  which criterion is decisive for an SOC class claim (joint signature vs
  single Vuong test), (iv) adding an explicit falsifier statement for
  Phase 2\textquotesingle s SOC verdict, and (v) adding the historical
  lognormal-vs-power-law econophysics literature context. An
  econophysics-savvy domain expert should still review the framing
  before submission. A separate decision about whether to issue a
  published correction note for the standalone arxiv-02 paper is the
  author\textquotesingle s editorial call.
\item
  \textbf{Sibling co-submission --- analysis written, author decides at
  submission (2026-05-24).} A pro/con analysis with a CC recommendation
  is in
  \texttt{docs/sessions/C1-v0.2-sibling-submission-decision-2026-05-24.md}.
  \textbf{CC recommendation: post C1 first; hold the thirteen-system
  sibling for 6--8 weeks, then post separately citing
  C1\textquotesingle s arXiv ID.} Rationale: epistemic dependency runs
  sibling → C1; provenance reconciliation cleaner in sequence; sibling
  has more failure modes (8+ Phases vs C1\textquotesingle s 4 real
  systems + null). Author makes the call at submission time.
\item
  \textbf{Domain-expert review --- internal proxy review written
  (2026-05-24); real review still required.} A three-reviewer-hat
  internal pre-submission review (seismology / econophysics /
  neuroscience) is in
  \texttt{docs/sessions/C1-v0.2-internal-review-2026-05-24.md}.
  \textbf{Summary: 9 P0 issues, 9 P1, 6 P2; the 3 hardest referee
  questions have draft defensible answers but the author should sign off
  on each before submission.} 5 of 9 P0 issues are pure editing fixes CC
  can apply directly in v0.3; 4 of 9 require re-running source-paper
  analyses (each ≤ 2 hours of focused work). The proxy review does
  \textbf{not} replace a real domain-expert pass; it is calibrated to a
  PRE / Chaos / Physica A referee level, not to BSSA / J. Finance / J.
  Neurosci. specialist referees. \emph{v0.4 addendum: §3.6 should
  additionally receive a complexity-science / statistical-mechanics
  domain-expert review on (a) the cross-domain scatter threshold,
  (b) the descriptor-vs-mechanism cluster, and (c) the MERGE
  recommendation for crackling-noise-universality.}
\item
  \textbf{Wave 3 follow-up plan (NEW in v0.4).} v0.4 §3.6 carries 3
  honest limitations (synthetic anchors for 11 classes; single-session
  verdicts; pre-reg bands occasionally over-specified). The Wave 3
  follow-up should: (i) replace synthetic anchors with real-data
  equivalents (priority on
  \texttt{adverse\_selection\_unraveling} comms-side BERTopic NLP,
  \texttt{gardner\_collins\_toggle\_v1} Anetzberger 2009 raw flow,
  \texttt{delay\_differential\_debt} PREDICTS / NOAA-ENSO / NSF-permafrost
  real data, \texttt{leaky\_integrate\_fire} SOEP real data,
  \texttt{anderson\_localization} published Anderson conductance
  datasets); (ii) add a second sub-agent verdict per class; (iii) revise
  pre-reg bands per the schelling-style mechanism-vs-magnitude
  distinction. Estimated effort: 10 sub-agents × \textasciitilde 60 min
  each = 10 wall-clock hours under 30-agent parallel mode.
\item
  \textbf{Taxonomy diagram update (NEW in v0.4).} The §3.6.7 textual
  spec should be rendered as
  \texttt{figures/taxonomy-v0.4.png} (PNG + SVG) before submission. The
  diagram should show: Layer 1 (10 mechanism nodes) + Layer 0
  (6 descriptor nodes) + Layer 2 (5+2 candidate nodes) + 5 SPLIT edges +
  1 MERGE edge + the cross-domain scatter threshold as a dashed
  horizontal screen. Estimated effort: 1 sub-agent × \textasciitilde 30
  min.
\end{enumerate}

\begin{center}\rule{0.5\linewidth}{0.5pt}\end{center}

\textbf{2026-05-24 v0.2 session deltas summary (CC closure work,
preserved from v0.2):}

\begin{itemize}
\tightlist
\item
  All 6 checklist items received CC-level closure work in the 2026-05-24
  v0.2 session.
\item
  3 supporting documents created in \texttt{docs/sessions/}:

  \begin{itemize}
  \tightlist
  \item
    \texttt{C1-v0.2-phase2-lognormal-revised-2026-05-24.md} (item 4)
  \item
    \texttt{C1-v0.2-sibling-submission-decision-2026-05-24.md} (item 5)
  \item
    \texttt{C1-v0.2-internal-review-2026-05-24.md} (item 6)
  \end{itemize}
\item
  1 git annotated tag created (item 2): \texttt{soc-pipeline-v0.1.0} at
  commit \texttt{4169928a}. Not pushed pending author authorization.
\item
  The C1 v0.2 draft itself was modified in-place for items 2, 3, and 4
  (inline edits to §2, §3.2, §6.1, References, Appendix A, and the
  Pre-submission checklist).
\item
  Items 1 and 5 still require author action before any submission (new
  Zenodo deposit; sibling co-submission decision).
\end{itemize}

\textbf{2026-05-24 v0.3 session deltas summary (P0 closure work):}

\begin{itemize}
\tightlist
\item
  All 9 P0 issues from the internal pre-submission review closed in
  v0.3:

  \begin{itemize}
  \tightlist
  \item
    4 pure-editing P0s (P0-E1, P0-E2, P0-N1, P0-N3) inlined into §3.2,
    §3.4, §6.1, §6.5, Table 1, References.
  \item
    5 re-run P0s (P0-S1, P0-S2, P0-S3, P0-E3, P0-N2) executed against
    existing source data; headline numbers in §3.1 and §3.4; full audit
    in Appendix A3.
  \end{itemize}
\item
  Re-run scripts at \texttt{/tmp/c1\_v03\_reruns.py} and
  \texttt{/tmp/c1\_v03\_p0\_n2\_per\_unit.py} (to be promoted to
  repository before submission tag --- see Appendix A3 reproducibility
  note).
\item
  0 P0 issues DEFERRED in v0.3.
\item
  3 new references added (LeBaron 2001, Malevergne et al. 2005,
  Pisarenko \& Sornette 2006; refs 27a/27b/27c).
\item
  Pre-submission checklist items 1 and 5 still require author action;
  item 6 (domain-expert review) now has its P0 issues closed at the CC
  layer but a real domain-expert pass remains required before
  submission.
\end{itemize}

\textbf{P0 disposition table (v0.3 final).}

{\def\LTcaptype{none} % do not increment counter
\begin{longtable}[]{@{}llll@{}}
\toprule\noalign{}
\# & P0 & Disposition & Where inlined \\
\midrule\noalign{}
\endhead
\bottomrule\noalign{}
\endlastfoot
1 & P0-S1 declustering robustness & RE-RUN & §3.1, Appendix A3 \\
2 & P0-S2 M\_c=4.45 FMD audit & RE-RUN & §3.1, Appendix A3 \\
3 & P0-S3 magnitude-type homogeneity & RE-RUN & §3.1, Appendix A3 \\
4 & P0-E1 daily-index scope & EDIT & §3.2 ("scope qualification"
para) \\
5 & P0-E2 lognormal econophysics refs & EDIT & §6.1 + Refs
27a/27b/27c \\
6 & P0-E3 Phase 2 Omori slope-zero & RE-COMPUTE & §3.2 (slope-zero
significance), Appendix A3 \\
7 & P0-N1 single-session caveat & EDIT & Table 1 footnote, §6.5 \\
8 & P0-N2 per-unit vs pooled & RE-RUN & §3.4 (pooled-spike vs per-unit),
Appendix A3 \\
9 & P0-N3 γ ≈ 1.10 vs γ\_MF = 2 & EDIT & §3.4 (scaling-relation γ
paragraph) \\
\end{longtable}
}

\begin{center}\rule{0.5\linewidth}{0.5pt}\end{center}

\textbf{2026-05-25 v0.4 session deltas summary (CC closure work):}

\begin{itemize}
\tightlist
\item
  v0.4 §3.6 "Completing the taxonomy" --- NEW section, including
  Table 2, the cross-domain scatter threshold §3.6.3, the 4-cluster
  cleanup §3.6.4, the two surprises §3.6.5, the synthetic-anchor
  honest-limit §3.6.6, and the textual taxonomy-figure spec §3.6.7. All
  18 verdicts cite their \texttt{v04-<class>-report.md} and
  \texttt{v4/validation/<class>/} artefacts.
\item
  Abstract rewritten to reflect 27 → 45+ SOC validation systems, KB
  main 4,888 entries + 445 pending-merge additions (merge ceiling
  5,333), the 18-class verdict matrix, the Layer-0 descriptor cluster,
  and the Spence-signal / OZ-Lorentzian surprises.
\item
  Table 1 preserved as v0.3 five-system core; Table 2 (v0.4 §3.6.2)
  added as the 18-class taxonomy-completion table.
\item
  Discussion §5 adds one paragraph on the methodological-classifier
  result.
\item
  Limitations §6 adds subsection 6.7 carrying §3.6\textquotesingle s
  three honest limitations.
\item
  References 46--50 added: in-repo verdict report set, C4
  mechanism-vs-descriptor paper, Halford 1992, Stumpf--Porter 2012,
  Cohen--Erez--ben-Avraham--Havlin 2000.
\item
  Appendix A4 added: per-class reproducibility map.
\item
  Pre-submission checklist items 7 + 8 added: Wave 3 follow-up plan,
  taxonomy diagram update.
\end{itemize}

\begin{center}\rule{0.5\linewidth}{0.5pt}\end{center}

\emph{End of draft v0.4. Status: 草稿待人审. Generated 2026-05-25 from
the v0.3 baseline + 17 v0.4 verdict reports + 1 leaky-integrate-fire
verdict.md. Number of pre-submission items: 8 (6 preserved from v0.3,
2 new in v0.4). Not yet reviewed by a real domain expert.}

\end{document}
